\bookStart{Speeches of Webthrithner}[Vaf·þrúðnis mǫ́l]
\setBookCode{Vafthrudnismal}

\begin{flushright}%
\textbf{Dating} \parencite{Sapp2022}: C10th (0.894)

\textbf{Meter:} \Ljodahattr%
\end{flushright}%

\section{Introduction}

The \textbf{Speeches of Webthrithner} (abbrev.~\Vafthrudnismal) describes a wisdom-contest between the god Weden and the ettin Webthrithner.

\subsection{Textual preservation and transmission}

\Vafthrudnismal\ is preserved in full only in \Regius\ (foll. 7v/9--8v/31), where it follows \textlink{Havamal} and is followed by \textlink{Grimnismal}.  The latter part of the poem, from st.~20 onwards, is also preserved in \AM\ (fol. 3r/1--3v/22).  Since the beginning of the poem has been lost in \AM, we cannot tell what text preceded it, but, as in \Regius, it is followed by \textlink{Grimnismal}.

\Vafthrudnismal\ was clearly available to Snorre.  Sts.~18, 30--31, 35, 37, 41, 45, 47, and 51 are cited in \Gylfaginning, and the lost form of the poem underlying \Gylfaginning\ may be called the \GylfMS\ redaction.  In the critical edition, \GylfMS\ will refer to agreement between all four mss.~of \Gylfaginning\ (\RegiusProse\Trajectinus\Wormianus\Upsaliensis).  It may be noted that st.~31 has lost its second helming (ll.~3--4) in both \Regius\ and \AM, which points towards a shared textual relationship between \Regius\AM\ to the exclusion of \GylfMS.

\subsection{Meter and stanza}

\Vafthrudnismal\ is composed in \Ljodahattr\ and consists of 55 stanzas.  Four lines may be considered defective and are irrestorable without conjectural emendation: 5/3, 34/1, 38/3 (see below), and 54/3.  In addition, st.~40 is heavily emended.

The normal stanzaic configuration \emph{abc abc} appears in 50 out of 55 sts.~(1--26, 28--37, 39--41, 44--54).  The singular configuration \emph{abc} (27; see note) in the middle of the question-and-answer stanzas should undoubtedly be considered a defective variant of this configuration.

Four variant configurations are found.  \emph{abc abcc} (42) and \emph{abcc abc} (43) are both variants of \emph{abc abc}.  They appear at an important transition in the poem, as Weden asks his last numbered question, introducing his final six questions about the end times.
\emph{abc abc abc} (55) only appears once, likewise at a charged location, namely in the final st., as Webthrithner admits defeat.

A stranger variant is \emph{abc ab abc} (38).  This configuration may be taken as a mixing of \Ljodahattr\ and \Malahattr, but it is also possible that a \emph{c}-verse has been lost following the central long-line \emph{ab}.  The fact that the long-line does not alliterate further hints at potential corruption, but it is clearly semantically crucial and cannot be removed from the stanza.

\subsection{Narrative structure}

The poem essentially consists of a riddle contest between the god Weden and the ettin Webthrithner.  Far from being a loose collection of mythic lore, it has a tight structure and logical plan throughout.  The whole may be divided into four sections: the first is the prologue, where Weden takes counsel from his wife Frie and sets out for Webthrithner’s hall (sts.~1--10).  The remaining three sections form the contest and consist of alternating stanzas where one side asks and the other answers.  They are distinguished from each other by means of repeated refrains in the question stanzas and consist of Webthrithner’s four unnumbered questions (11--19), Weden’s twelve numbered questions (20--43), and Weden’s six unnumbered questions about the end times (44--55).

The following table illustrates the refrains; for stanza 40 see note there:

\begin{small}
\begin{center}
\begin{tabular}{|c c|}
  \hline
  11--17 & \emph{Sęg mér/þat, Gang·ráðr, \hld\ alls ȧ golfi vill / þïns of fręista frama} \\
  20--42 & \emph{Sęg þat (it) \emph{N}(:a) \hld\ ... / ... Vaf·þrúðnir vitir} \\
  44--54 & \emph{Fjǫlð ek fór, \hld\ fjǫlð fręistaða’k, / fjǫlð of ręynda’k ręgin} \\
  \hline
\end{tabular}
\end{center}
\end{small}

Something ought to be said about the numerology of the questions.  It is hardly a coincidence that Weden asks exactly eighteen questions, this being a multiple of the sacred number nine.  It is especially notable that another Odinic list, Leed-tally (sts.~147--165) of \textlink{Havamal}, also has eighteen items---even more so, since the eighteenth spell there, like the eighteenth question here, is a mystery known only to Weden himself.  This circumstance reveals the care put into the structure of \Vafthrudnismal.

\subsection{Summary}

The first scene of the poem assumes that the god Weden and his wife, Frie, are sitting together.  Weden asks Frie for counsel, as he is curious about the ancient wisdom which the ettin Webthrithner might possess (1).  Frie expresses worry, since she considers Webthrithner stronger than all other ettins (2), but Weden says that he has travelled far and wide and wants to know what Webthrithner’s hall is like (3).  Frie wishes him good luck against the ettin (4), and he departs to challenge Webthrithner’s \emph{orð-spęki} ‘word-wisdom’ (5).  The god arrives at the ettin’s hall and introduces himself (6); Webthrithner promptly declares that his guest will not come out of the hall unless he be wiser than he (7).  Weden introduces himself as Gangred, saying that he has come from afar in need of Webthrithner’s hospitality (8).  The ettin invites this Gangred to sit down (9), and the guest in turn utters a gnomic stanza (10) not unlike those of the first section of \textlink{Havamal}.

Webthrithner begins by asking four mythological questions, each answered by Gangred in turn. The questions concern the horse that pulls the day (11--12), the horse that pulls the night (13--14), the river which divides the gods and ettins (15--16), and the plain where \inx[P]{Surt} and the gods will fight (17--18).

Webthrithner calls the guest learned and invites him to sit.  He declares that the loser of the contest must give his head (19).  The roles are now reversed, and Gangred poses twelve numbered questions to the ettin.  The questions can be divided into three groups of four based on their subjects.  Gangred first asks about the origins of cosmological phenomena: the earth and heaven (20--21), sun and moon (22--23), day and night (24--25), and winter and summer (26--27).  His next four questions concern the earliest ettins: the ettin \inx[P]{Earyelmer} (28--29), Earyelmer’s ancestry (30--31), how he reproduced without a partner (32--33), and what Webthrithner himself first remembers (34--35).  The last four numbered questions are more sundry.  Weden asks about the origin of the wind (36--37) and the god \inx[P]{Nearth} (38--39), about Walhall and the Oneharriers (40--41), and at last a meta-question: where Webthrithner learned such great wisdom (42--43).

The refrain and theme of the questions now change, as Gangred asks six unnumbered questions concerning the end times, all beginning with the same refrain.  The disguised god asks which humans will survive the \inx[P]{Fimblewinter} (44--45), how the sun can rise after Fenrer has destroyed it (46--47), about some obscure maidens (48--49; see note), which Eese will survive after the flame of Surt goes out (50--51), and how Weden will die (52--53).  Finally, he asks the unknowable question: what did Weden speak in the ear of Balder before he was burned on the pyre (54)?

Webthrithner at last understands the identity of his challenger, since only Weden himself could know the answer to that question.  He laconically accepts his imminent death and the futility of his own wisdom; the poem ends with his admission that Weden will always be the wisest (55).

\subsection{Dating}

\Vafthrudnismal\ displays archaic formal characteristics.  The poem has an especially high frequency of the archaic particle \emph{of}, appearing at a rate of 0.438 occurrences per 10 unique lines \parencite[70]{Sapp2022}.  In addition, it displays alliteration between \emph{v-} and old \emph{vr-} (st.~53/2; \cite[37]{Haukur2016}), a highly archaic trait indicative of composition before ca.~1000 \ce.  The pre-Christian subject matter, which shows no sign of antiquarian composition, and which was likely not fully understood even by Snorre (see st.~35), also suggests an early date.  In conclusion, \Vafthrudnismal\ was likely composed in the 9th or 10th century---at any rate during the pre-Christian period.

\newpage

\section{Text}

\subsection{Vaf·þrúðnis mǫ́l (‘The Speeches of Webthrithner’)}

\bvg\bva\mssnote{\Regius~7v/9}%
\speakernote{[Óðinn:]}%
\llap{„}Ráð mér nú \alst{F}rigg \hld\ alls mik \alst{f}ara tíðir &
\ind at \alst{v}itja \alst{V}af·þrúðnis; &
\edtext{\alst{f}or·vitni mikla \hld\ kveð’k mér ȧ \alst{f}ornum stǫfum &
\ind við þann inn \alst{a}l-svinna \alst{jǫ}tun.}{\lemma{for·vitni \dots\ jǫtun. ‘Great \dots\ ettin.’}\Bfootnote{I.e., “I am very curious to learn his ancient words of wisdom.”  Cf.~st.~55.}}“\eva

\bvb\speakernoteb{[Weden:]}%
“{\huge C}\textsc{ounsel me now, \inx[P]{Frie},} as I long to journey \\
\ind to visit \inx[P]{Webthrithner}. \\
Great curiosity I have for the ancient staves \\
\ind of that all-wise \inx[P]{Ettins}[ettin].”\evb\evg


\bvg\bva\mssnote{\Regius~7v/12}%
\speakernote{[Frigg:]}%
\llap{„}\alst{H}ęima lętja \hld\ mynda’k \edtrans{\alst{H}ęrja-fǫðr}{Father of Hosts}{\Bfootnote{A name for Weden.  For the second element, see \textlink{Voluspa}[1]/3 n.}} &
\ind ï \alst{g}ǫrðum \alst{g}oða; &
því’t \alst{ę}ngi \alst{jǫ}tun \hld\ hugða’k \alst{ja}fn-ramman &
\ind sęm \alst{V}af·þrúðni \alst{v}esa.“\eva

\bvb\speakernoteb{[Frie:]}%
“At home would I keep the Father of Hosts \\
\ind in the yards of the Gods, \\
for no ettin have I judged to be \\
\ind as strong as Webthrithner.”\evb\evg


\bvg\bva\mssnote{\Regius~7v/13}%
\speakernote{[Óðinn:]}%
\llap{„}\alst{F}jǫlð ek \alst{f}ór, \hld\ \alst{f}jǫlð \alst{f}ręistaða’k, &
\ind fjǫlð \edtext{\emph{of}}{\Afootnote{emend.~(cf.~st.~44/2); om.~\Regius}} \alst{r}ęynda’k \alst{r}ęgin; &
hitt \alst{v}il’k \alst{v}ita, \hld\ hvé \alst{V}af·þrúðnis &
\ind \alst{s}ala-kynni \alst{s}éi.“\eva

\bvb\speakernoteb{[Weden:]}%
“Much I journeyed, much I tried, \\
\ind much I tested the \inx[P]{Reins}! \\
One more thing I wish to know: how Webthrithner’s \\
\ind halls may be.”\evb\evg


\bvg\bva\mssnote{\Regius~7v/15}%
\speakernote{[Frigg:]}%
\llap{„}\alst{H}ęill þú farir, \hld\ \alst{h}ęill þú aptr komir, &
\ind hęill þú ȧ \alst{s}innum \alst{s}éir; &
\alst{ǿ}ði þér dugi \hld\ hvar’s \edtext{skalt}{\Afootnote{add.~\emph{†or†} \Regius}}, \edtrans{\alst{A}lda-fǫðr}{Father of Mankind}{\Bfootnote{A rare poetic name for Weden, also found in st.~53/1, but not otherwise used in Eddic poetry.  Among genuine compositions, the name is only found elsewhere in the Scaldic Bragi \emph{Þórr} 1/1 (in \Skp\ 3), although it also appears in an antiquarian poetic list of names for Weden, Þul \emph{Óðins} 1/2 (in \Skp\ 3).}}, &\nobreak
\ind \alst{o}rðum mę́la \alst{jǫ}tun.“\eva

\bvb\speakernoteb{[Frie:]}%
“Hale mayst thou journey; hale mayst thou come back; \\
\ind hale mayst thou be on thy paths! \\
May thy wisdom avail thee where thou, Father of Mankind, \\
\ind with words shalt greet the ettin!”\evb\evg


\bvg\bva\mssnote{\Regius~7v/17}%
\alst{F}ór þȧ Óðinn \hld\ at \alst{f}ręista orð-spęki &
\ind þęss ins \alst{a}l-svinna \alst{jǫ}tuns; &
at hǫllu hann kom, \hld\ \edtext{\emph{es}}{\Afootnote{emend.; \emph{ok} \Regius}} átti \edtrans{Ïms}{Ime’s}{\Bfootnote{Gen.~sg.~of \emph{Ïmr}, an unknown ettin.  The name is probably corrupt, since alliteration on \emph{h-} is required by the strongly stressed \emph{hǫllu} in the a-verse.  \textcite{FinnurEdda} emends to \emph{Hymis} ‘Hymer’s’.}} faðir; &
\ind \alst{i}nn gekk \edtrans{\alst{Y}ggr}{Ug}{\Bfootnote{‘The Terrifier’, another name for Weden.}} þegar.\eva

\bvb Then journeyed Weden to test the word-wisdom \\
\ind of that all-wise ettin. \\
He came to the hall which Ime’s father \ken*{= Webthrithner} owned; \\
\ind \inx[P]{Ug} \name{= Weden} went soon inside.\evb\evg


\bvg\bva\mssnote{\Regius~7v/18}%
\speakernote{[Óðinn:]}%
\llap{„}\alst{H}ęill þú nú, Vaf·þrúðnir, \hld\ nú em’k ï \alst{h}ǫll kominn &
\ind ȧ þik \alst{s}jalfan \alst{s}éa; &
hitt vil’k \alst{f}yrst vita, \hld\ ef \alst{f}róðr séir &
\ind eða \alst{a}l-sviðr, \alst{jǫ}tunn.“\eva

\bvb\speakernoteb{[Weden:]}%
“Hale be thou now, Webthrithner! Now I am come into the hall \\
\ind to see thy very self, \\
yet I wish first to know if thou be learned \\
\ind or all-wise, ettin!”\evb\evg


\bvg\bva\mssnote{\Regius~7v/20}%
\speakernote{[Vaf·þrúðnir:]}%
\llap{„}Hvat ’s þat \alst{m}anna, \hld\ es ï \alst{m}ïnum sal &
\ind verp⸗umk \alst{o}rði \alst{ȧ}? &
\alst{ú}t þú né kømr \hld\ \alst{ȯ}rum hǫllum frȧ, &
\ind nema þú inn \alst{s}notrari \alst{s}éir.“\eva

\bvb\speakernoteb{[Webthrithner:]}%
“What sort of man is this who in \emph{my} hall \\
\ind throws his word at me? \\
Out wilt thou not come from our halls \\
\ind unless thou be the cleverer man!”\evb\evg


\bvg\bva\mssnote{\Regius~7v/22}%
\speakernote{[Óðinn:]}%
\llap{„}\edtext{\alst{G}a\emph{ng}·ráðr}{\Afootnote{\emph{Gagn·ráðr} \Regius}\Bfootnote{‘One who decides to walk’.  The emendation is supported by the b-verse, where \emph{gǫngu} (dat.~sg.~of \emph{ganga} ‘walking, walk’) is probably tying back to the name.
\emph{Gang·ráðr} is also found in Þul \emph{Óðins} 3/1 (\Skp\ 3, Gurevich ed.).
The variant \emph{Gagn·ráðr} would mean ‘Victory-advisor’ and is not attested elsewhere.  Note that Gurevich is incorrect when she states in her note to Þul \emph{Óðins} 3/1 that \Gylfaginning~5:12 records the variant \emph{Gagn·ráðr}.  The relevant source mss., \RegiusProse\Trajectinus, only have the abbreviation \emph{Þá spurði G.} ‘Then G.~asked’, which clearly does not point one way or the other.}} hęiti’k, \hld\ nú em’k af \alst{g}ǫngu kominn, &
\ind \alst{þ}yrstr til \alst{þ}ïnna sala; &
\alst{l}aðar þurfi \hld\ hęf’k \alst{l}ęngi farit &
\ind ok þïnna \alst{a}nd-fanga, \alst{jǫ}tunn.“\eva

\bvb\speakernoteb{[Weden:]}%
“\inx[P]{Gangred} I am called; now I am come from walking, \\
\ind thirsty, to thy halls. \\
In need of a welcome I’ve journeyed for long, \\
\ind and of thy reception, ettin!”\evb\evg


\bvg\bva\mssnote{\Regius~7v/24}%
\speakernote{[Vaf·þrúðnir:]}%
\llap{„}Hví þú þȧ, \alst{G}ang·ráðr, \hld\ mę́li⸗sk af \alst{g}olfi fyrir? &
\ind far þú ï \alst{s}ess ï \alst{s}al; &
þȧ skal \alst{f}ręista, \hld\ hvárr \alst{f}lęira viti, &
\ind \alst{g}ęstr eða \edtrans{inn \alst{g}amli þulr}{the old thyle}{\Bfootnote{Webthrithner himself, the thyle being the lorekeeper whose purpose it was to recite the old wisdom poems.  See Encyclopedia: \inx[C]{thyle}.}}.“\eva

\bvb\speakernoteb{[Webthrithner:]}%
“Why then, Gangred, dost thou speak from off the floor ahead? \\
\ind Take a seat in the hall! \\
Then it shall be tried which of the two might know more: \\
\ind the guest, or the old \inx[C]{thyle}.”\evb\evg


\bvg\bva\mssnote{\Regius~7v/26}%
\speakernote{[Óðinn:]}%
\llap{„}\alst{Ȯ}·auðigr maðr, \hld\ es til \alst{au}ðigs kømr, &
\ind \edtrans{mę́li \alst{þ}arft eða \alst{þ}ęgi}{ought to speak the needful or shut up}{\Bfootnote{Formulaic, this line occurs identically in \textlink{Havamal}[19].}}; &
\alst{o}fr-mę́lgi mikil \hld\ hygg’k at \alst{i}lla geti &
\ind hvęim’s við \edtrans{\alst{k}ald-rifjaðan}{cold-ribbed}{\Bfootnote{Cold-hearted, cunning.}} \alst{k}ømr.“\eva

\bvb\speakernoteb{[Weden:]}%
“An unwealthy man who to a wealthy one comes \\
\ind ought to speak the needful or shut up. \\
Great over-speaking I think will bring ill \\
\ind for whomever comes upon a cold-ribbed man.”\evb\evg


\bvg\bva\mssnote{\Regius~7v/28}%
\Ballnote{Sts.~11--14 deal with the two horses that pull the day and night through the heavens.  Their origin is discussed in \Gylfaginning~10:1--10 (the whole chapter): \emph{Nǫrvi eða Narfi hét jǫtunn, er byggði í Jǫtun-heimum.  Hann átti dóttur, er Nótt hét.  Hon var svǫrt ok døkk, sem hon átti ę́tt til.  Hon var gipt þeim manni, er Nagl·fari hét; þeira sonr hét Auðr.  Því nę́st var hon gipt þeim, er Ánarr hét; Jǫrð hét þeira dóttir.  Síðast átti hana Dellingr, ok var hann ása ę́ttar; var þeira sonr Dagr; var hann ljóss ok fagr eftir faðerni sínu.  Þá tók Al·fǫðr Nótt ok Dag, son hennar, ok gaf þeim tvá hesta ok tvę́r kerrur ok sendi þau upp á himin, at þau skulu ríða á hverjum tveim dǿgrum um·hverfis jǫrð’ina.  Ríðr Nótt fyrri þeim hesti, er kallaðr er Hrím·faxi, ok at morgni hverjum dǫggvir hann jǫrð’ina af mél-dropum sínum.  Sá hestr, er Dagr á, heitir Skin·faxi, ok lýsir allt lopt ok jǫrð’ina af faxi hans.} ‘Narrow or Narve was the name of an ettin who dwelled in the Ettinhomes.  He had a daughter whose name was Night.  She was black and dark like her lineage bespoke.  She was married off to the man whose name was Nailfarer.  Their son was Ead.  Thereafter she was married off to the one whose name was Oner.  Earth was the name of their daughter.  Lastly Delling got her, and he was of the lineage of the Eese.  Their son was Day; he was bright and fair after his father’s side.  Then Allfather took Night, and Day her son, and gave them two horses and two chariots and sent them up into the heaven that they should ride every two half-days a full circle around the earth.  Night rides in front on the horse which is called Rimefax, and every morning he bedews the earth with the foam from his bit.  The horse owned by Day is named Shinefax, and all the air and the earth is lit up by his mane.’

They are mentioned again in \Skaldskaparmal\ 72, which deals with legendary horses: \emph{Ár·vakr ok Al·sviðr draga sól’ina, sem fyrr er ritat. Hrím·faxi eða Fjǫr·svartnir draga nótt’ina. Skin·faxi eða Glaðr fylgja degi’num.} ‘Yorewaker and Allswith pull the sun as is written earlier.  Rimefax or Ferrowswartner pulls the night.  Shinefax or Glad follows the day.’   For Yorewaker and Allswith see \textlink{Grimnismal}[38].  The alternative names Ferrowswartner (‘Life-darkener’) and Glad probably derive from now-lost poetic sources.

The motif of horses, especially mounted on chariots, pulling the heavenly bodies or day and night is very ancient, and is attested as far back as the Bronze Age.  It is common in the hymns of \Rigveda\ and Greek texts, and is represented archeologically by the beautiful Trundholm Sun Chariot from Denmark of the Nordic Bronze Age, the cultural and ethnic ancestors of the Germanic peoples.  See Figures \ref{fig:trundholm-front} and \ref{fig:trundholm-back}, and cf.~\textlink{Grimnismal}[38] n.}%
\speakernote{[Vaf·þrúðnir:]}%
\llap{„}Sęg mér, \alst{G}ang·ráðr, \hld\ alls ȧ \alst{g}olfi vill &
\ind \edtrans{þïns of \alst{f}ręista \alst{f}rama}{test thy furtherance}{\Bfootnote{I.e., “try your luck, see how far you get”.  Formulaic; cf.~\textlink{Havamal}[2]/4.}}, &
hvé \alst{h}ęstr \alst{h}ęitir, \hld\ sá’s \alst{h}vęrjan dręgr &
\ind \alst{d}ag of \alst{d}rótt-mǫgu.“\eva

\bvb\speakernoteb{[Webthrithner:]}%
“Tell me, Gangred, since on the floor thou wilt \\
\ind test thy furtherance, \\
what the horse is called which pulls every \\
\ind day over the lads of the folk \ken{men}.”\evb\evg


\bvg\bva\mssnote{\Regius~7v/30}%
\speakernote{[Óðinn:]}%
\llap{„}\alst{Sk}in·faxi hęitir, \hld\ es inn \alst{sk}íra dręgr &
\ind \alst{d}ag of \alst{d}rótt-mǫgu; &
\alst{h}ęsta batstr \hld\ þykkir hann með \edtext{\emph{\alst{H}}ręið-gotum}{\Afootnote{metr. emend.; \emph{‘reið-gotom’} \Regius}\lemma{\emph{H}ręið-gotum ‘Reth-Gots’}\Bfootnote{An old name for the Eastern Gots around the Black Sea, also found in \textlink{Widsith}[12][57] and Ög 136 Rök.  The origin of the first element is unclear, but it appears to have a tribal significance (cf.~the word \emph{Hréð-cyning} ‘Reth-King’ in \textlink{Widsith}[1][7]).  They are apparently mentioned here due to their connection with the east, the direction of the dawn.  There may also be a pun present, since \emph{goti} can mean both ‘Got’ and ‘horse’.}}; &
\ind ęy lýsir \alst{m}ǫn af \alst{m}ari.“\eva

\bvb\speakernoteb{[Weden:]}%
“\inx[P]{Shinefax} is he called who pulls the pure \\
\ind day over the lads of the folk. \\
The best of horses he seems among the \inx[P]{Reth-Gots}; \\
\ind for ever that stallion’s mane lights up.”\evb\evg

\begin{figure}[b]
\centering
\includegraphics[width=\textwidth]{Solvognen-front}
\caption{The gilt (day) side of the Trundholm sun chariot.  Nordic Bronze Age, ca. 1400 \bce.  © Nationalmuseet. \href{https://creativecommons.org/licenses/by-sa/3.0/deed.en}{CC BY-SA 3.0}.  \url{https://commons.wikimedia.org/wiki/File:Solvognen_Do_2010_1277.tif}}
\label{fig:trundholm-front}
\end{figure}

\bvg\bva\mssnote{\Regius~7v/32}%
\speakernote{[Vaf·þrúðnir:]}%
\llap{„}Sęg þat, \alst{G}ang·ráðr, \hld\ alls ȧ \alst{g}olfi vill &
\ind þïns of \alst{f}ręista \alst{f}rama, &
hvé \alst{jó}r hęitir, \hld\ sá’s \alst{au}stan dręgr &
\ind \alst{n}ǫ́tt of \edtrans{\alst{n}ýt ręgin}{the bountiful Reins}{\Bfootnote{The Gods.  \emph{nýtr} ‘bountiful’ is one of several positive adjectives that can stand before \emph{ręgin}, and recurs in the same position in \textlink{Vafthrudnismal}[14]/2, 25/3 below.  Cf.~\emph{holl ręgin} ‘the hold Reins’ (\textlink{Lokasenna}[4]/3), \emph{fróð ręgin} ‘the learned Reins’ (\textlink{Vafthrudnismal}[26]/4), and \emph{blíð ręgin} ‘the blithe Reins’ (\textlink{Grimnismal}[6]/1, 38/3, 42/1; sarcastically in \textlink{Lokasenna}[32]/3).}}.“\eva

\bvb\speakernoteb{[Webthrithner:]}%
\emph{“Tell this, Gangred, since on the floor thou wilt \\
\ind test thy furtherance,} \\
what the steed is called which from the east does pull \\
\ind night over the bountiful \inx[P]{Reins}.”\evb\evg

\begin{figure}[b]
\centering
\includegraphics[width=\textwidth]{Solvognen-back}
\caption{The ungilt (night) side of the Trundholm sun chariot.  © Nationalmuseet. \href{https://creativecommons.org/licenses/by-sa/3.0/deed.en}{CC BY-SA 3.0}.  \url{https://commons.wikimedia.org/wiki/File:Solvognen_Do_2010_1278.tif}}
\label{fig:trundholm-back}
\end{figure}

\bvg\bva\mssnote{\Regius~7v/33}%
\speakernote{[Óðinn:]}%
\llap{„}\edtrans{\alst{H}rím·faxi}{Rimefax}{\Bfootnote{The one with a frozen mane.}} \alst{h}ęitir, \hld\ es \alst{h}vęrja dręgr &
\ind \alst{n}ǫ́tt \edtext{of}{\Afootnote{emend.; \emph{ok} \Regius}} \alst{n}ýt ręgin; &
\alst{m}ėl-dropa \hld\ fęllir hann \alst{m}orgin hvęrjan; &
\ind \edtrans{þaðan kømr \alst{d}ǫgg of \alst{d}ala}{thence comes dew through the dales}{\Bfootnote{For another explanation of the origin of dew, see \textlink{Voluspa}[18].}}.“\eva

\bvb\speakernoteb{[Weden:]}%
“\inx[P]{Rimefax} is he called who pulls each \\
\ind night over the bountiful Reins. \\
Foam from his bit he lets fall each morning; \\
\ind thence comes dew through the dales.”\evb\evg


\bvg\bva\mssnote{\Regius~8r/1}%
\speakernote{[Vaf·þrúðnir:]}%
\llap{„}Sęg þat, \alst{G}ang·ráðr, \hld\ alls ȧ \alst{g}olfi vill &
\ind þïns of \alst{f}ręista \alst{f}rama, &
hvé \alst{ǫ́} hęitir, \hld\ sú’s dęilir með \alst{jǫ}tna sonum &
\ind \alst{g}rund, ok með \alst{g}oðum.“\eva

\bvb\speakernoteb{[Webthrithner:]}%
\emph{“Tell this, Gangred, since on the floor thou wilt \\
\ind test thy furtherance,} \\
what the river is called which divides the land \\
\ind between the sons of Ettins and the Gods.”\evb\evg


\bvg\bva\mssnote{\Regius~8r/2}%
\speakernote{[Óðinn:]}%
\llap{„}\edtrans{\alst{Í}fing}{Iving}{\Bfootnote{The border river is not known by this name from any other source, not even \Gylfaginning, which otherwise tends to relay even the most obscure lore.}} hęitir \alst{ǫ́}, \hld\ es dęilir með \alst{jǫ}tna sonum &
\ind \alst{g}rund, ok með \alst{g}oðum; &
\alst{o}pin rinna \hld\ hǫ̇n skal umb \alst{a}ldr-daga; &
\ind \edtrans{verðr⸗at \alst{í}ss ȧ \alst{ǫ́}u}{there forms no ice on that river}{\Bfootnote{For ice would enable the Ettins to cross over into the lands of the Gods.}}.“\eva

\bvb\speakernoteb{[Weden:]}%
“\inx[P]{Iving} is the river called which divides the land \\
\ind between the sons of Ettins and the Gods. \\
Open shall it flow through its days of life; \\
\ind there forms no ice on that river.”\evb\evg


\bvg\bva\mssnote{\Regius~8r/3}%
\speakernote{[Vaf·þrúðnir:]}%
\llap{„}Sęg þat, \alst{G}ang·ráðr, \hld\ alls ȧ \alst{g}olfi vill &
\ind þïns of \alst{f}ręista \alst{f}rama, &
hvé \alst{v}ǫllr hęitir, \hld\ es finna⸗sk \alst{v}ígi at &
\ind \alst{S}urtr ok in \alst{s}vǫ́su goð.“\eva

\bvb\speakernoteb{[Webthrithner:]}%
\emph{“Tell this, Gangred, since on the floor thou wilt \\
\ind test thy furtherance,} \\
what the plain is called where they find each other at war, \\
\ind \inx[P]{Surt} and the beloved Gods.”\evb\evg


\bvg\bva\mssnote{\Regius~8r/4, \RegiusProse~17r/24, \\
\Trajectinus~TODO, \Wormianus~TODO}%
\Ballnote{For the fight between Surt and the Gods cf.~\textlink{Voluspa}[48]--53, \textlink{Fafnismal}[14]--15.  The st.~is explained or paraphrased in \Gylfaginning~51:25: \emph{Mú·spells megir sǿkja fram á þann vǫll er Víg·ríðr heitir.} ‘The lads of Muspell make their attack on the plain which is called Wighride.’ and again in \Gylfaginning~51:28--29: \emph{En Mú·spells synir hafa einir sér fylking, ok er sú bjǫrt mjǫk.  Vǫllr’inn Víg·ríðr er hundrað rasta víðr á hvern veg.} ‘But the sons of Muspell have a single formation, and it is very bright.  The plain Wighride is a hundred rests wide in each direction.’  The stanza itself is cited in \Gylfaginning~51:51, the very end of that chapter, following the citation of \textlink{Voluspa}[45]--54, introduced by the words \emph{Hér segir enn svá} ‘Here it further says so’.  Note that it is omitted from \Upsaliensis, one of the four \Gylfaginning\ mss.}%
\speakernote{Óðinn:}%
\llap{„}\alst{V}íg·ríðr hęitir \alst{v}ǫllr, \hld\ es finna⸗sk \alst{v}ígi at &
\ind \alst{S}urtr ok in \alst{s}vǫ́su goð; &
\alst{h}undrað \edtrans{rasta}{rests}{\Bfootnote{An old measurement of distance, ca.~1850 metres (\CV, s.v.~\emph{rǫst}).}} \hld\ hann ’s ȧ \edtext{\alst{h}vęrjan}{\Afootnote{\emph{hvęrn} \Trajectinus}} veg; &
\ind sá ’s þęim \edtext{\alst{v}ǫllr}{\Afootnote{add.~\emph{of} \Trajectinus}} \alst{v}itaðr.“\eva

\bvb\speakernoteb{Weden:}%
“\inx[P]{Wighride} is the plain called where they find each other at war, \\
\ind Surt and the beloved Gods. \\
A hundred \inx[C]{rest}[rests] it reaches in each direction; \\
\ind for them that plain is marked out.”\evb\evg


\bvg\bva\mssnote{\Regius~8r/6}%
\speakernote{Vaf·þrúðnir:}%
\llap{„}\edtrans{\alst{F}róðr est nú gęstr, \hld\ \alst{f}ar á bękk jǫtuns}{Learned art thou now, guest; take the ettin’s bench}{\Bfootnote{I.e., “Thou hast proven thy learning.”  Webthrithner acknowledges his guest’s wisdom and repeats his invitation for the god to sit down.  Weden accepts the offer, and the contest continues; now it is Weden’s turn to ask, Webthrithner’s to answer.}}, &
\ind ok mę́lum⸗k ï \alst{s}essi \alst{s}aman; &
\alst{h}ǫfði vęðja \hld\ vit skulum \alst{h}ǫllu ï &
\ind \alst{g}ęstr, of \alst{g}oð-spęki.“\eva

\bvb\speakernoteb{Webthrithner:}%
“Learned art thou now, guest; take the ettin’s bench \\
\ind and let us speak in the seat together! \\
Wager a head shall we two in the hall, \\
\ind O guest, over god-wisdom!”\evb\evg

\sectionline

\subsection{Capitulum (‘New Chapter’)}

\bvg\bva\mssnote{\Regius~8r/9, \AM~3r/1}%
\speakernote{Óðinn:}%
\llap{„}\edtext{Sęg}{\Afootnote{\emph{S} is a bold capital in \Regius}} þat it \alst{ęi}na, \hld\ ef þïtt \edtext{\alst{ǿ}ði}{\lemma{ǿði}\Afootnote{The first word on fol.~3r of \AM \emph{;} from this point we have the poem in both manuscripts.}} dugir &
\ind \edtext{ok}{\Afootnote{\emph{eða} \AM}} þú \alst{V}af·þrúðnir \alst{v}itir, &
hvaðan \edtext{\alst{jǫ}rð of kom, \hld\ eða \alst{u}pp-himinn}{\lemma{jǫrð \dots\ eða upp-himinn ‘Earth \dots\ or Up-heaven’}\Bfootnote{An old Common Germanic formulaic merism, referring to the whole universe.  See Index of formulae: \inx[F]{Earth and Up-heaven}.}} &
\ind \alst{f}yrst, inn \alst{f}róði jǫtunn.“\eva

\bvb\speakernoteb{Weden:}%
“{\huge T}\textsc{ell this one} if thy wisdom avails \\
\ind and thou, Webthrithner, mightst know \\
whence Earth came or \inx[P]{Up-heaven}, \\
\ind first, O learned ettin.”\evb\evg


\bvg\bva\mssnote{\Regius~8r/10, \AM~3r/2}%
\speakernote{Vaf·þrúðnir kvað:}%
\Ballnote{The gods sacrificed Yimer and created the world from his body, as told more fully in \textlink{Grimnismal}[41]--42 and alluded to in \textlink{Voluspa}[3]--4.  For Yimer, cf.~sts.~30--33 below; for the deeper cosmological and religious significance of this myth see notes to \textlink{Voluspa}[3]--4 and also to \textlink{Grimnismal}[43]. --- The whole st.~bears very close resemblance to \textlink{Grimnismal}[41]; ll.~1 and 4 here are identical to ll.~1--2 there, and ll.~2 and 3a here are clearly related to ll.~3a and 4 there.  Still, the sts.~are distinct enough that the one cannot be a direct scribal copy of the other, and the relationship is more likely to be oral.  Both have probably been composed in the same West Norwegian milieu, in turn deriving from an older Common Germanic tradition (cf.~\textlink{Wessobrunn} under Poetry on Christian Subjects).}%
\llap{„}Ór \alst{Y}mis holdi \hld\ vas \alst{jǫ}rð of skǫpuð, &
\ind ęn ór \alst{b}ęinum \alst{b}jǫrg, &
\edtrans{\alst{h}iminn ór \alst{h}ausi}{the heaven from the skull}{\Bfootnote{The heavens are understood as a dome, a view common to many ancient peoples.  This also fits well with the floating clouds being Yimer’s brains, as told in \textlink{Grimnismal}[42].}} \hld\ ins \alst{h}rím-kalda jǫtuns, &
\ind ęn \edtrans{ór \edtrans{\alst{s}vęita}{blood}{\Bfootnote{In poetry \emph{svęiti} ‘sweat’ almost always means ‘blood’. This is shared with OE \emph{swât}, as seen e.g. in \Beowulf\ 1286a: \emph{sweord \emph{swâte} fâh} ‘sword stained with “sweat”’, 2689b--2690: \emph{hé ge-blódegod wearð / sâwul-dríore; \hld\ \emph{swât} ýðum wéoll.} ‘he was bloodied in soul-gore; the “sweat” gushed in waves’.}} \alst{s}ę́r}{from his blood the sea}{\Bfootnote{According to \Gylfaginning~7:3--4, the slaying of Yimer produced so much blood that it drowned the whole race of Rime-Thurses save one; see st.~35 below.
For the method of slaying cf.~\Sonatorrek\ 3/3: \emph{jǫtuns hals \hld\ undir þjóta} ‘the neck-wounds of the ettin \ken{seas} roar’, which attests that Yimer was slain by decapitation, the typical way of wasting beasts of sacrifice (so e.g. \textlink{Hymiskvida}[15]).
That this is not a mere literary construct is proven by the excavation of the Wiking Age site Hove-steads (\emph{Hofstaðir}) on Iceland, where bulls were seasonally slain in what was undoubtedly ritual sacrifice: “The most likely reconstruction from the forensics of the skulls requires at least a two-person team, one of whom struck the animal between the eyes (effectively killing it and certainly stunning it into momentary immobility) while the second swung a fairly broad-bladed axe at the neck or base of the skull for a beheading stroke.” (\cite[23]{Lucas2007}).  After the slaying (and presumed feasting on the meat), the heads of the bulls were displayed for a prolonged period of time.
\citeauthor{Lucas2007} further note that swift beheading was not the usual manner of slaughtering animals on Iceland, and even has practical downsides compared to a slower cutting of the throat, e.g., splintered bones and damage to the cutting blade.  On the other hand, the violent cutting and following flow of blood would have great dramatic effect, and, what \citeauthor{Lucas2007} do not mention, clearly reenact the slaying of Yimer: the separation of the skull (heaven) from the body (earth), and the great flow of blood (sea-water) from the neck-wound, lastly the burial of the consumed body in the earth, and the display of the skull on high to symbolize the heaven. ---

The ancients believed that the earth was surrounded on all sides by an outer sea (ON \emph{út-haf}); cf.~\textlink{Hymiskvida}[5]/1 n., 22/4 n.}}.“%NOTE: this is a hack for parencite which for whatever reason doesn't work inside a Bfootnote with multiple authors. reledmac...
\eva

\bvb\speakernoteb{Quoth Webthrithner:}%
“From \inx[P]{Yimer}’s flesh was the earth created, \\
\ind but from his bones the mountains; \\
the heaven from the skull of that rime-cold ettin, \\
\ind but from his blood the sea.”\evb\evg


\bvg\bva\mssnote{\Regius~8r/12, \AM~3r/3}%
\speakernote{Óðinn kvað:}%
\Ballnote{For the origin of the Sun and Moon cf.~\textlink{Grimnismal}[38].}%
\llap{„}Sęg þat \alst{a}nnat, \hld\ \edtext{ef þïtt \alst{ǿ}ði dugir &
\ind ok þú \alst{V}af·þrúðnir \alst{v}itir}{\lemma{ef \dots\ vitir}\Afootnote{abbrev.~as \emph{‘e. þ. ę. d. ⁊ þ. v.’} \Regius \emph{;} abbrev.~as \emph{‘e. þ. e. d.’} \AM}}, &
hvaðan \alst{M}ȧni of kom, \hld\ svá’t fęrr \alst{m}ęnn yfir, &
\ind eða \alst{S}ól it \alst{s}ama.“\eva

\bvb\speakernoteb{Quoth Weden:}%
\emph{“Tell this} other one, \emph{if thy wisdom avails,} \\
\ind \emph{and thou, Webthrithner, mightst know} \\
whence the Moon came which journeys over men, \\
\ind or the Sun likewise.”\evb\evg


\bvg\bva\mssnote{\Regius~8r/13, \AM~3r/4}%
\speakernote{Vaf·þrúðnir kvað:}%
\Ballnote{This st.~clearly serves as a source of \Gylfaginning~11:1--3: \emph{Þá mę́lti Gangleri: „Hversu stýrir hann gang sólar eða tungls?“ Hárr segir: „Sá maðr er nefndr Mundil·fari, er átti tvau bǫrn.  Þau vǫ́ru svá fǫgr ok fríð, at hann kallaði son sinn Mána, en dóttur sína Sól ok gipti hana þeim manni, er Glenr hét. [\dots]} ‘Then spoke Gangler: “How does He [viz. Allfather] rule the track of the sun or the moon?” High says: “There is a man named Mundlefarer who had two children.  They were so fair and beautiful that he called his son ‘Moon’, but his daughter ‘Sun’, and he married her off to the man who was called Glen. [\dots]’  Cf.~\textlink{Grimnismal}[38] n., which cites \Gylfaginning~11:4--7.}%
\llap{„}\edtrans{\alst{M}undil·fǿri}{Mundlefarer}{\Bfootnote{See Index for the etymology, which likens the cosmos to a Wiking Age flour-mill turned by a handle.  His name is attested in \Gylfaginning~11:2 as \emph{Mundil·fari}, which underlies the Anglish rendition Mundlefarer.}} hęitir, \hld\ hann ’s \alst{M}ȧna faðir &
\ind ok svá \alst{S}ólar it \alst{s}ama; &
\alst{h}imin \alst{h}verfa \hld\ þau skulu \alst{h}vęrjan dag &
\ind \edtrans{\alst{ǫ}ldum at \edtext{\alst{á}r-tali}{\Afootnote{corr.~in same hand from \emph{ald-tali} \AM}}}{for mankind’s tally of years}{\Bfootnote{According to \textlink{Voluspa}[6] the Gods gave names to night, the moon-phases, morning, midday, afternoon, and evening \emph{ǫ́rum at tęlja} ‘the years for to tally’. --- Numerous examples of the chronological reckoning of the Heathen Icelanders are found in Are’s Book of Icelanders and in the Book of Land-Claims (\Landnamabok).  Both texts anchor the years to the Christian Common Era, but this is likely the result of later Christian learning.  The counts were most likely originally based on the reigns of kings, of which many examples are found in those two books.}}.“\eva

\bvb\speakernoteb{Quoth Webthrithner:}%
“\inx[P]{Mundlefarer} he is called---he is the father of the Moon, \\
\ind and so of the Sun likewise. \\
Turn round heaven shall they every day, \\
\ind for mankind’s tally of years.”\evb\evg


\bvg\bva\mssnote{\Regius~8r/15, \AM~3r/6}%
\speakernote{Óðinn kvað:}%
\llap{„}\alst{S}ęg þat it þriðja, \hld\ alls þik \alst{s}vinnan kveða &
\ind \edtext{\edtext{ok}{\Afootnote{\emph{eða} \AM}} þú \alst{V}af·þrúðnir \alst{v}itir}{\Afootnote{abbrev.~as \emph{‘⁊ þ. v. v.’} \Regius}}, &
hvaðan \alst{D}agr of kom, \hld\ sá’s fęrr \alst{d}rótt yfir, &
\ind eða \edtrans{\alst{N}ǫ́tt með \alst{n}iðum}{Night with the moon-phases}{\Bfootnote{Also found in \textlink{Voluspa}[6]/3.}}.“\eva

\bvb\speakernoteb{Quoth Weden:}%
\emph{“Tell this} third, \emph{since they call thee wise,} \\
\ind \emph{and thou, Webthrithner, mightst know}, \\
whence the Day came which journeys over the folk, \\
\ind or the Night with the moon-phases.”\evb\evg


\bvg\bva\mssnote{\Regius~8r/17, \AM~3r/8}%
\speakernote{Vaf·þrúðnir kvað:}%
\llap{„}\edtrans{\alst{D}ęllingr}{Delling}{\Bfootnote{A derivative of the adjective \emph{dallr} ‘(possibly) shining, brilliant’ (cf.~Old English \emph{deal} ‘proud, exalted’) which appears in two other ON mythological names, viz. \emph{Hęim·dallr} (the god \inx[P]{Homedall}) and \emph{Mar·dǫll} (Meredall, an epithet of the goddess \inx[P]{Frow}.  Also in these two instances the root appears to have connotations of light and brilliance; for Meredall see \textlink{Voluspa}[24] n., \textlink{Thrymskvida}[8] n.,
and for Homedall see \textlink{Grimnismal}[13], \textlink{Thrymskvida}[13]/3 n., \Gylfaginning~27 (ed.~under \textlink{EddicFragments}[F3][III]).}} hęitir, \hld\ hann ’s \alst{D}ags faðir, &
\ind ęn \alst{N}ǫ́tt vas \edtrans{\alst{N}ǫrvi}{Narrow}{\Bfootnote{Dat.~sg.~of \emph{Nǫrr}, to be identified with the \emph{Nǫrvi} (or \emph{Narfi}) of \Gylfaginning~10:1 (see st.~11 n.).  His name may be a kenning for the Gap of Ginnings or primæval chaos (cf.~\textlink{Voluspa}[3]/4 n.) and thereby equivalent to the Greek \textgreek{χάος} ‘gap, Chaos’, father of \textgreek{Νύξ} ‘Night’ (\Theogony\ 123).}} borin; &
\edtrans{\alst{n}ý ok \alst{n}ið}{The waxing and waning}{\Bfootnote{The two halves of the moon cycle, by which months (’moon-ths’) were reckoned.  The phrase also occurs in \Gutasaga\ 2: \emph{um ný ok niðar}, \Gulatingslog\ 57: \emph{um ný hit nę́sta ok niðar}, and \Gylfaginning~11:8: \emph{Máni stýrir gǫngu tungls ok rę́ðr nýjum ok niðum} ‘Moon steers the track of the lunar body and rules over the waxings and wanings’.  For the Norse-Germanic calendar see notes to \textlink{Voluspa}[6].}} \hld\ skópu \alst{n}ýt ręgin &
\ind \alst{ǫ}ldum at \alst{á}r-tali.“\eva

\bvb\speakernoteb{Quoth Webthrithner:}%
“\inx[P]{Delling} he is called; he is the father of \inx[P]{Day}, \\
\ind but \inx[P]{Night} was born to \inx[P]{Narrow}. \\
The waxing and waning the bountiful Reins created \\
\ind for mankind’s tally of years.”\evb\evg


\bvg\bva\mssnote{\Regius~8r/18, \AM~3r/9}%
\speakernote{Óðinn kvað:}%
\llap{„}Sęg þat it \alst{f}jórða, \hld\ \edtext{alls þik \edtrans{\alst{f}\emph{róðan}}{learned}{\Bfootnote{It seems indisputable that the word hiding behind the abbreviated \emph{f.} of the mss.~is the m.~acc.~sg.~of the adj.~\emph{fróðr}, but as it is not written out explicitly it must be marked by italics.}} kveða, &
\ind ok þú \alst{V}af·þrúðnir \alst{v}itir}{\lemma{alls \dots\ vitir}\Afootnote{abbrev.~as \emph{‘a. þ. f.’} \Regius\AM}}, &
hvaðan \alst{v}etr of kom \hld\ eða \alst{v}armt sumar &
\ind \alst{f}yrst \edtext{með}{\Afootnote{\emph{of} \AM}} \alst{f}róð ręgin.“\eva

\bvb\speakernoteb{Quoth Weden:}%
\emph{“Tell this} fourth, \emph{since they call thee learned,} \\
\ind \emph{and thou, Webthrithner, mightst know} \\
whence the Winter came or the warm Summer, \\
\ind first, amidst the learned Reins.”\evb\evg


\bvg\bva\mssnote{\Regius~8r/20, \AM~3r/10}%
\Ballnote{This st.~is paraphrased in \Gylfaginning~19:3--4: \emph{Svásuðr heitir sá, er faðir Sumars er, ok er hann svá sę́l-lífr, svá at af hans heiti er þat kallat svásligt, er blítt er. En faðir Vetrar er ýmist kallaðr Vind·lóni eða Vind·svalr.  Hann er Vásaðar son, ok vǫ́ru þeir ǫ́ttungar grimmir ok sval-brjóstaðir, ok hefir Vetr þeira skap-lyndi.} ‘Sosuth is the name of he who is the father of Summer, and he lives so joyously that from his name that which is glad is called ‘delightful’.  But the father of Winter is variously called Windloon or Windswale.  He is the son of Wosuth, and those kinsmen were cruel and cool-breasted, and Winter has their temperament.’}%
\speakernote{Vaf·þrúðnir:}%
\llap{„}\alst{V}ind·svalr hęitir, \hld\ hann’s \alst{V}etrar faðir, &
\ind ęn \alst{S}vǫ́suðr \alst{S}umars.“ &
\edtext{[...]}{\Bfootnote{The second half of the st.~seems to be missing.  Its contents are unknown, but may be preserved in the \Gylfaginning~paraphrase.  No lacuna is indicated in the mss.}}\eva

\bvb\speakernoteb{Webthrithner:}%
“\inx[P]{Windswale} is he called; he is the father of the \inx[P]{Winter}; \\
\ind but \inx[P]{Sosuth} is the \inx[P]{Summer}’s.”\evb\evg


\bvg\bva\mssnote{\Regius~8r/21, \AM~3r/11}%
\speakernote{Óðinn kvað:}%
\llap{„}Sęg þat it \alst{f}imta, \hld\ \edtext{alls þik \alst{f}\emph{róðan} kveða, &
\ind ok þú \alst{V}af·þrúðnir \alst{v}itir}{\lemma{alls \dots\ vitir}\Afootnote{abbrev.~as \emph{‘a. þ. f.’} \Regius \emph{;} abbrev.~as \emph{‘a. þ. f. k.’} \AM}}, &
\edtext{hvęrr \alst{ȧ}sa \alst{ę}ldstr \hld\ eða \alst{Y}mis niðja &
\ind \alst{y}rði ï \alst{á}r-daga.}{\lemma{hvęrr \dots\ ár-daga ‘who \dots\ days of yore.’}\Bfootnote{I.e., “Which was the very first being?”  Cf.~the question on the cryptic C9th Malt Stone (DR NOR1988;5): \textbf{huaʀisi : alistiąsa}, perhaps \emph{Hvaʀ es inn ęldsti ȧsa?} ‘Who is the eldest of the Eese?’}}“\eva

\bvb\speakernoteb{Quoth Weden:}%
\emph{“Tell this} fifth, \emph{since they call thee learned,} \\
\ind \emph{and thou, Webthrithner, mightst know} \\
who, oldest of the \inx[P]{Eese} or of Yimer’s kinsmen \ken{ettins}, \\
\ind arose in days of yore.”\evb\evg


\bvg\bva\mssnote{\Regius~8r/22, \AM~3r/12}%
\speakernote{Vaf·þrúðnir kvað:}%
\llap{„}\alst{Ø}r·ófi vetra \hld\ áðr vę́ri \alst{jǫ}rð of skǫpuð, &
\ind þȧ \edtext{vas}{\Afootnote{om.~\AM}} \alst{B}er·gęlmir \alst{b}orinn; &
\alst{Þ}rúð·gęlmir \hld\ vas \alst{þ}ess faðir, &
\ind ęn \alst{Au}r·gęlmir \alst{a}fi.“\eva

\bvb\speakernoteb{Quoth Webthrithner:}%
“Uncountable winters before the Earth was created, \\
\ind then was \inx[P]{Bareyelmer} born. \\
\inx[P]{Thrithyelmer} was that one’s father, \\
\ind and \inx[P]{Earyelmer} the grandfather.”\evb\evg


\bvg\bva\mssnote{\Regius~8r/23, \AM~3r/14, \\
\RegiusProse~2v/17, \Trajectinus~3r/28, \\
\Wormianus~TODO}%NOTE: This stanza is not found in U.
\Ballnote{In \Gylfaginning~5:12, Snorre cites \Vafthrudnismal\ 30/3--4 and the full 31 in support of his lengthy and embellished creation narrative, introduced by the words \emph{En hér segir svá Vaf·þrúðnir jǫtunn:} ‘but here the ettin Webthrithner says so’, though it is not certain that his vision is what the poet of \Vafthrudnismal\ had in mind.}%
\speakernote{Óðinn kvað:}%
\llap{„}\edtext{\edtext{\alst{S}ęg þat it}{\Afootnote{abbrev.~as \emph{‘S. þ. þ. h.’} (for *\emph{S. þ. h.}) \AM}} \alst{s}étta, \hld\ \edtext{alls þik \alst{s}vinnan kveða, &
\ind ok þú \alst{V}af·þrúðnir \alst{v}itir,}{\lemma{alls \dots\ vitir}\Afootnote{abbrev.~as \emph{‘a. þ. s. q.’} \Regius \emph{;} abbrev.~as \emph{‘a. þ. f. k.’} (influenced by sts.~26, 28) \AM}}}{\lemma{Sęg \dots\ vitir ‘Tell \dots\ know’}\Afootnote{om. \GylfMS}} &
hvaðan \alst{Au}r·gęlmir kom \hld\ með \alst{jǫ}tna sonum &
\ind \alst{f}yrst, inn \alst{f}róði jǫtunn.“\eva

\bvb\speakernoteb{Quoth Weden:}%
\emph{“Tell this} sixth, \emph{since they call thee wise,} \\
\ind \emph{and thou, Webthrithner, mightst know} \\
whence Earyelmer came amidst the sons of ettins, \\
\ind first, O learned ettin.”\evb\evg


\bvg\bva\mssnote{\Regius~8r/25, \AM~3r/15, \\
\RegiusProse~2v/17, \Trajectinus~3r/29, \\
\Wormianus~TODO, \Upsaliensis~TODO}%
\Ballnote{Over æons the splashing droplets from the primæval venomous rivers, the Ilewaves, combined until they formed a sentient being.  This was Earyelmer, whom \Gylfaginning~5:9 (most likely correctly) identifies with \inx[P]{Yimer}, for whom cf.~\textlink{Voluspa}[3]/1 n.}%
\speakernote{Vaf·þrúðnir kvað:}%
\llap{„}\edtext{\alst{Ó}r}{\Afootnote{\emph{Þȧ’s ór} \RegiusProse}} \edtrans{\alst{É}li-vǫ́gum}{Ilewaves}{\Bfootnote{The wild ettin-ish rivers near the edge of the world; cf.~\textlink{Hymiskvida}[5]/1.

The Ilewaves (\emph{Éli-vágar}) are likely a reflex of the primæval salty Waters found in many West Eurasian mythologies, including Genesis 1:1--3 and \Rigveda\ 10.129.  Of those two foundational religious sources the latter agrees more closely with the present stanza and certainly holds the more archaic conception; for the Jewish Waters (Hebrew \texthebrew{תְּהוֹם} \emph{tᵉhôm}) and their Mesopotamian and Egyptian predecessors cf.~\textcite{Waschke2006}.

Whereas we in the Jewish narrative find a \emph{creation out of nothing}---at the very beginning of time Elohim’s spirit is on the Waters and He makes the light shine over them---we instead find in the two Indo-European texts \Vafthrudnismal\ and \Rigveda\ a \emph{spontaneous emergence} of a single primæval entity long before the Gods are born, here from the violent splashing of venom, in \Rigveda\ 10.129.3 from “the power of heat” (OI \textoi{tápasas maɦinā́}).
This entity in turn asexually begets sexual beings---here through rubbing his limbs together (\textlink{Vafthrudnismal}[33]), in \Rigveda\ 10.129.4 by giving rise to “desire” (\textoi{kā́ma}) which serves as the “primal seed of thought” (\textoi{mánasas rḗtas pratʰamám}).  It is from these beings that the world is populated by the raw matter out of which the Gods eventually create (or \emph{shape}) the World as we know it through the Creative Act of the First Sacrifice (see st.~21 above, \textlink{Voluspa}[3]--4, 6--7, \textlink{Grimnismal}[41]--43).

It should be noted that the Ilewaves also correspond to the Greek Earth-bounding river \textgreek{Ὠκεανός}, especially in its earlier conception; cf.~\textlink{Hymiskvida}[5]/1 n.}} \hld\ stukku \edtrans{\alst{ęi}tr-dropar}{droplets of venom}{\Bfootnote{This \emph{ęitr} ‘venom’ of the Waves is probably to be understood as referring to extreme saltiness.  Yimer’s blood gives rise to the salty sea; and the stoniness of his bones, perhaps originally thought of as salt-stone, is attested for both his skull (Heaven, the stoney Firmament; \textlink{Harbardsljod}[15]/2) and his body (the Earth and Bedrock; st.~21/2, \textlink{Grimnismal}[41]/3).  To this day, something very salty is said to be \emph{ettersalt} in modern Swedish.}}, &
\ind \edtrans{svá}{so}{\Afootnote{\emph{ok} ‘and’ \GylfMS}} \edtext{\alst{ó}x und’s \edtext{ór varð}{\Afootnote{rev.~\Trajectinus\Wormianus}} \alst{jǫ}tunn}{\Afootnote{\emph{vǫxtr vinds ok varð jǫtunn ór} ‘the growth of the wind and there became an ettin from it’ (for *\emph{óx unds varð jǫtunn ór}) \Upsaliensis}}; &
\edtext{\edtext{þar}{\Afootnote{\emph{þę́r} \Upsaliensis \emph{;} add. \emph{eru} \RegiusProse}} \edtext{\alst{ȯ}rar}{\Afootnote{\emph{†onar†} \Trajectinus, \emph{ęinar} \Upsaliensis}} \alst{ę́}ttir \hld\ \edtext{kómu}%
{\Afootnote{\emph{komnar} \RegiusProse\emph{;} \emph{koma} \Upsaliensis}} \edtext{\alst{a}llar}{\Afootnote{om.~\Upsaliensis}} saman; &
\ind \edtext{því’s \edtrans{þat}{it}{\Bfootnote{The ettin race.}} \edtext{\alst{ę́}}{\Afootnote{om.~\Trajectinus\Wormianus}} \alst{a}lt \edtext{til}{\Afootnote{\emph{eða} \Trajectinus \emph{;} om.~\Wormianus}} \alst{a}talt}{\Afootnote{om. \Upsaliensis}}.“}{\lemma{þar \dots\ atalt ‘Thereof \dots\ fierce’}\Afootnote{om. \Regius\AM.}}\eva

\bvb\speakernoteb{Quoth Webthrithner:}%
“Out of the \inx[P]{Ilewaves} splashed droplets of venom; \\
\ind so it grew until there became an ettin from it. \\
Thereof our lineages all came together; \\
\ind therefore it is for ever all too fierce.”\evb\evg


\bvg\bva\mssnote{\Regius~8r/26, \AM~3r/16}%
\speakernote{Óðinn kvað:}%
\llap{„}\edtext{\alst{S}ęg þat it}{\Afootnote{abbrev.~as \emph{‘S. þ. h.’} \AM}} \alst{s}jaunda, \hld\ \edtext{alls þik \alst{s}vinnan kveða, &
\ind ok þú \alst{V}af·þrúðnir \alst{v}itir}{\lemma{alls \dots\ vitir}\Afootnote{abbrev.~as \emph{‘a. þ. s. q.’} \Regius \emph{;} abbrev.~as \emph{‘a. þ. s. k.’} \AM}}, &
hvé sá \alst{b}ǫrn gat \hld\ inn \edtrans{\alst{b}aldni}{stubborn}{\Afootnote{so \AM\emph{;} \emph{aldni} ‘olden’ \Regius}} jǫtunn, &
\ind es hann hafði⸗t \alst{g}ýgjar \alst{g}aman.“\eva

\bvb\speakernoteb{Quoth Weden:}%
\emph{“Tell this} seventh, \emph{since they call thee wise,} \\
\ind \emph{and thou, Webthrithner, mightst know} \\
how he begot children, that stubborn ettin, \\
\ind when he knew no troll-woman’s pleasure.”\evb\evg


\bvg\bva\mssnote{\Regius~8r/27, \AM~3r/17}%
\Ballnote{This stanza is paraphrased in \Gylfaginning~5:16--18: \emph{Ok svá er sagt, at þá er hann svaf, fekk hann sveita.  Þá óx undir vinstri hendi hónum maðr ok kona, ok annarr fótr hans gat son við ǫðrum.  En þaðan af kómu ę́ttir.  Þat eru hrím-þursar.} ‘And so it is said that when he slept he began to sweat.  Then grew within his left hand a man and a woman, and one foot of his begat a son by the other.  But thereof come lineages.  Those are Rime-Thurses.’}%
\speakernote{Vaf·þrúðnir kvað:}%
\llap{„}\edtext{Und \alst{h}ęndi vaxa \hld\ kvǫ́ðu \alst{h}rím-þursi &
\ind \alst{m}ęy ok \alst{m}ǫg saman; &
\alst{f}ótr við \alst{f}ǿti}{\lemma{Und hęndi \dots\ fótr við fǿti ‘Under the arm \dots\ Foot against foot’}\Bfootnote{The image is masturbatory and monstrous; for its significance cf.~st.~31 n.}} \hld\ \edtext{gat}{\Afootnote{om.~\AM}} ins \alst{f}róða jǫtuns &
\ind \alst{s}ex-hǫfðaðan \alst{s}on.“\eva

\bvb\speakernoteb{Quoth Webthrithner:}%
“Under the arm of the \inx[P]{Rime-Thurses}[rime-thurse], they said, did grow \\
\ind a maiden and a lad together. \\
Foot against foot begat for the learned ettin \\
\ind a six-headed son.”\evb\evg


\bvg\bva\mssnote{\Regius~8r/29, \AM~3r/18}%
\speakernote{Óðinn kvað:}%
\llap{„}\edtrans{\edtext{\alst{S}ęg þat it}{\Afootnote{abbrev.~as \emph{‘S. þ. h.’} \AM}} ǫ́ttunda, \hld\ alls þik f\emph{róðan} kveða}{Tell this eighth, since they call thee learned}{\Bfootnote{This line lacks the required alliteration, which has probably been lost in manuscript transmission due to the repetitive nature of the refrain.  Alliteration may be obtained by replacing \emph{alls þik fróðan kveða} with \emph{ef þïtt ǿði dugir} from sts.~20 and 22, or (less favourably) with \emph{alls þik svinnan kveða} from 24, where the alliteration in the a-verse does not fall on the ordinal.}}\edtext{, &
\ind ok þú \alst{V}af·þrúðnir \alst{v}itir}{\lemma{alls \dots\ vitir}\Afootnote{abbrev.~as \emph{‘a. þ. f.’} \Regius \emph{;} abbrev.~as \emph{‘a. þ. f. k.’} \AM}}, &
\linenumannotation{a}hvat \alst{f}yrst \edtext{of}{\Afootnote{om.~\Regius}} mant \hld\ eða \alst{f}ręmst of vęitst, &
\ind þú est \alst{a}l-sviðr \alst{jǫ}tunn.“\eva

\bvb\speakernoteb{Quoth Weden:}%
\emph{“Tell this} eighth, \emph{since they call thee learned,} \\
\ind \emph{and thou, Webthrithner, mightst know} \\
what first thing thou recallest or earliest knowest--- \\
\ind thou art all-wise, ettin!”\evb\evg


\bvg\bva\mssnote{\Regius~8r/30, \AM~3r/19, \\
\RegiusProse~3r/4, \Trajectinus~3v/20, \\
\Wormianus~TODO, \Upsaliensis~TODO}%
\Ballnote{This st.~is cited in \Gylfaginning~7:6, introduced by the words \emph{svá sem hér segir} ‘just as it says here’.

The preceding \Gylfaginning~7:2--5 interprets it in the following way: the sons of \inx[P]{Byre} (that is, \inx[P]{Weden}, \inx[P]{Will}, and \inx[P]{Wigh}) slew the ettin Yimer, but when he died so much blood flowed from his wounds that the whole race of Ettins was drowned, save for Bareyelmer and his household, who survived by getting up on his \emph{lúðr} (cf.~note to this word below).

The flood-myth in \Gylfaginning~7:2--5 is clearly related to the Great Flood or Deluge as known from many mythologies, most famously Genesis 6--9.  Although it may be a genuine Germanic reflex of the same ancient myth (cf.~\textlink{Voluspa}[4] n., 56/1 n.), it may also be Snorre’s own invention based on the Bible in an attempt to explain the present stanza, which may in fact have been as obscure to him as it is to us.}%
\speakernote{Vaf·þrúðnir kvað:}%
\llap{„}\alst{Ø}r·ófi vetra \hld\ áðr vę́ri \alst{jǫ}rð of skǫpuð, &
\ind þȧ vas \alst{B}er·gęlmir \alst{b}orinn; &
þat \edtext{ek}{\Afootnote{\emph{er} \Trajectinus\Wormianus}} \alst{f}yrst of man, \hld\ \edtext{es inn \alst{f}róði}{\Afootnote{\emph{at fróða} \Upsaliensis}} jǫtunn &
\ind \edtext{ȧ vas}{\Afootnote{so \AM\RegiusProse\Wormianus\Upsaliensis \emph{;} rev.~\Regius\Trajectinus}} \edtrans{\alst{l}úðr}{canoe}{\Bfootnote{In Old Norse prose \emph{lúðr} usually means ‘trumpet, blowing horn’, less commonly ‘flour-bin’, but the underlying sense seems to be ‘hollowed-out wood’, which is why it is presently translated as ‘canoe’.
As noted above Snorre seems to interpret it as a form of ark on which Bareyelmer found refuge from the flood, but considering the transitive nature of Bareyelmer’s being laid (\emph{of lagiðr}) upon it, the stanza may instead be interpreted as speaking of a ship burial, so that the first thing Webthrithner remembers is Bareyelmer’s funeral.}} of \alst{l}agiðr.“\eva

\bvb\speakernoteb{Quoth Webthrithner:}%
“Uncountable winters before the Earth was created, \\
\ind then was Bareyelmer born. \\
This I first remember, when the learned ettin \\
\ind on the canoe was laid.”\evb\evg


\bvg\bva\mssnote{\Regius~8r/32, \AM~3r/21}%
\Ballnote{Sts.~36--37 are paraphrased by \Gylfaginning~18:1--6: \emph{Þá mę́lti Gangleri: „Hvaðan kemr vindr? Hann er sterkr, svá at hann hrę́rir stór hǫf, ok hann ę́sir eld. En svá sterkr sem hann er, þá má eigi sjá hann; því er hann undar⸗liga skapaðr.“ Þá mę́lti Hárr: „Þat kann ek vel segja þér.  Á norðan-verðum himins enda sitr jǫtunn sá, er Hrę́·svelgr heitir.  Hann hefir arnar-ham, en er hann beinir flug, þá standa vindar undan vę́ngjum hans.} ‘Then Gangler spoke: “Whence comes the wind? It is so strong that it stirs great oceans and it fans flame.  But as strong as it is yet one cannot see it.  Therefore it is wondrously created.” Then High spoke: “That I can tell thee well.  At the northern end of heaven sits the ettin who is called Rawswallower.  He has an eagle’s hame, and when he takes to flight then the winds come from under his wings.”’

After this, st.~37 is cited in \Gylfaginning~18:7, introduced by the words \emph{Hér segir svá} ‘Here it says so’.}%
\speakernote{Óðinn kvað:}%
\llap{„}\edtext{\alst{S}ęg þat it}{\Afootnote{abbrev.~as \emph{‘S. þat h.’} \AM}} níunda, \hld\ \edtext{alls þik \alst{s}vinnan kveða, &
\ind ok þú \alst{V}af·þrúðnir \alst{v}itir}{\lemma{alls \dots\ vitir}\Afootnote{abbrev.~as \emph{‘a. þ. s. q. e.’} \Regius \emph{;} \emph{‘a. þ. s. k.’} \AM}}, &
hvaðan \alst{v}indr of kømr \hld\ svá’t fęrr \alst{v}ág yfir, &
\ind \edtrans{ę́ \edtext{męnn hann \alst{s}jalfan of \alst{s}éa}{\Afootnote{\emph{maðr of sjalfan hann sér} \AM}}}{always men see his very self}{\Bfootnote{Perhaps a reference to the sea which is never perfectly still, so that the wind is always seen on the waves.  It is also possible that a negative clitic \emph{-t} has been lost from the verb \emph{séa} ‘see’, in which case the line should be interpreted as ‘\emph{never} men see his very self’.  The loss of an original negative is supported by \Gylfaginning~18:3: \emph{má eigi sjá hann} ‘one cannot see it’.}}.“\eva

\bvb\speakernoteb{Quoth Weden:}%
\emph{“Tell this} ninth, \emph{since they call thee wise,} \\
\ind \emph{and thou, Webthrithner, mightst know} \\
whence the wind comes which fares o’er the wave--- \\
\ind always men see his very self.”\evb\evg


\bvg\bva\mssnote{\Regius~8r/34, \AM~3r/22, \\
\RegiusProse~5v/21, \Trajectinus~TODO, \\
\Wormianus~TODO, \Upsaliensis~TODO}%
\speakernote{Vaf·þrúðnir kvað:}%
\llap{„}\edtrans{\alst{H}rę́·svęlgr}{Rawswallower}{\Bfootnote{‘Devourer of Carrion’.}} \edtext{\alst{h}ęitir}{\Afootnote{add.~\emph{jǫtunn} \Upsaliensis}}, \hld\ \edtext{es}{\Afootnote{\emph{hann} \Upsaliensis}} sitr ȧ \alst{h}imins ęnda, &
\ind \alst{jǫ}tunn ï \alst{a}rnar ham; &
af hans \alst{v}ę̇ngjum \hld\ kveða \alst{v}ind \edtext{koma}{\Afootnote{\emph{standa} \Upsaliensis}} &
\ind \alst{a}lla męnn \alst{y}fir.“\eva

\bvb\speakernoteb{Quoth Webthrithner:}%
“\inx[P]{Rawswallower} is he called who sits at heaven’s end, \\
\ind an ettin in an eagle’s \inx[C]{hame}. \\
From his wings they say the wind comes \\
\ind over all men.”\evb\evg


\bvg\bva\mssnote{\Regius~8v/1, \AM~3r/24}%
\speakernote{Óðinn kvað:}%
\llap{„}\edtext{\alst{S}ęg þat it}{\Afootnote{abbrev.~as \emph{‘S. þat h.’} \AM}} \alst{t}íunda, \hld\ \edtext{alls}{\Afootnote{abbrev.~as \emph{‘a.’} \Regius\AM}} þú \edtrans{\alst{t}íva rǫk}{rakes of the Tews}{\Bfootnote{The fated events or narrations of the Gods.}} &
\ind ǫll \alst{V}af·þrúðnir \alst{v}itir, &
\edtext{hvaðan Njǫrðr of kom \hld\ með ȧsa sonum}{\Bfootnote{This line lacks both alliteration and a following c-verse, which may point towards textual corruption.  \textcite{FinnurEdda} opts to change the b-verse to *\emph{með niðjum ȧsa}, a perfectly sensible emendation, although the expected c-verse is still missing.  For that reason, I instead wish to propose the following conjectural restoration: \emph{hvaðan Njǫrðr of kom \hld} af Nóa-tu̇num / ár \emph{með ȧsa sonum} ‘whence Nearth of the Nowetons came of yore in the midst of the sons of the Eese’, where *\emph{af Nóa-tu̇num / ár} would have been omitted due to the visual similarity between *\emph{tu̇num} and \emph{sonum}.}}; &
\edtrans{\alst{h}ofum ok \alst{h}ǫrgum}{hoves and harrows}{\Bfootnote{A formulaic merism for cultic sites.  \emph{hǫrgr} ‘harrow’ is a cultic place, not the agricultural implement; cf.~\textlink{Voluspa}[7] n.  For Nearth in connection with harrows cf.~\textlink{Grimnismal}[16], where it is said that Nearth \emph{rę́ðr hǫ́-timbruðum hǫrgi} ‘rules the harrow timbered on high’.
Also of interest is \textlink{Lokasenna}[51], where the goddess Shede speaks of her \emph{vé ok vangar} ‘wighs and wongs’, two terms common in religious place names.  The underlying theological understanding seems to be that the god or goddess is physically present in his or her shrine as its ruler \parencite[238]{TurvillePetre1975}.}} \hld\ hann rę́ðr \edtrans{\alst{h}und-mǫrgum}{a hundred-many}{\Bfootnote{Probably a reference to the large number of cultic sites named after Nearth; among modern place names \textcite{Brink2007} counts 13 in Norway, 17 in Sweden, and 3 in Denmark, but those numbers must have been much greater in the pre-Christian period.}} &
\ind ok \edtext{varð⸗at}{\Afootnote{\emph{vas⸗at} \AM}} hann \alst{ǫ̇}sum \alst{a}linn.“\eva

\bvb\speakernoteb{Quoth Weden:}%
“Tell this tenth, since thou the rakes of the Tews \\
\ind all, Webthrithner, mightst know, \\
whence \inx[P]{Nearth} came in the midst of the sons of the \inx[P]{Eese}. \\
Of \inx[C]{hove}[hoves] and \inx[C]{harrow}[harrows] he rules a hundred-many, \\
\ind and he was not begotten by the Eese.”\evb\evg


\bvg\bva\mssnote{\Regius~8v/3, \AM~3r/26}%
\Ballnote{In later prose sources we learn that following the settlement between the Eese and the Wanes (i.e., the two tribes of Gods), Nearth, who belonged to the Wanes, was sent to the Eese as a hostage.  Note however that \emph{goð} ‘Gods’ in l.~2 here seems to refer exclusively to the Eese.  For this narrative cf.~\textlink{Voluspa}[22] n., \Gylfaginning~23:6 (edited in the pres. ed. as \textlink{EddicFragments}[F2][II]), and \YnglingaSaga\ 4.
\textlink{Lokasenna}[34]--35 mentions that Nearth was sent as a hostage to the Ettins, but that is certainly a different episode, although the shared motif of Nearth as a hostage is notable.}%
\speakernote{Vaf·þrúðnir kvað:}%
\llap{„}İ \alst{V}ana-hęimi \hld\ skópu hann \edtrans{\alst{v}ís ręgin}{the wise Reins}{\Bfootnote{ON \emph{ręgin} ‘the Reins, Powers’ is normally one of a number of terms used to refer to the gods as a collective \parencite[10--11]{Simek2010}, but here seems to refer specifically to the \inx[P]{Wanes} in opposition to the \inx[P]{Eese}.}} &
\ind ok sęldu at \alst{g}íslingu \alst{g}oðum, &
ï \edtrans{\alst{a}ldar rǫk}{the Rakes of the Age}{\Bfootnote{The \inx[P]{Rakes of the Reins}, the End Times.}} \hld\ \edtext{hann mun}{\Afootnote{rev.~\AM}} \alst{a}ptr koma &
\ind hęim með \alst{v}ísum \alst{v}ǫnum.“\eva

\bvb\speakernoteb{Quoth Webthrithner:}%
“In \inx[P]{Waneham} the wise \inx[P]{Reins} created him \\
\ind and sold him as a hostage to the Gods. \\
In the \inx[P]{Rakes of the Age} he will come back \\
\ind home amidst the wise \inx[P]{Wanes}.”\evb\evg


\bvg\bva\mssnote{\Regius~8v/5, \AM~3r/28}%
\Ballnote{This question-stanza is malformed in both \Regius\ and \AM\ and thus has to be partly reconstructed on the basis of st.~41.  The ms. preservation of \Vafthrudnismal\ 40--41 is as follows:

All four mss.~of \Gylfaginning~attest only st.~41 (the answer-stanza, cf.~note there) with no textual variants.

\Regius\ has two half-stanzas which are clearly a mix between the question and the answer: \emph{Segðv þat et .xi. hvar ytar tvnom i hꜹggvaz hverian dag. Ynom. Val þeir kíosa oc riþa vígi fra sitia meirr vm sáttir saman.}  Although this forms a complete stanza it stands out by lacking a refrain in the first two lines, something found in all other question-sts.~in \Vafthrudnismal\ (see Introduction).  The word \emph{Ynom} (where the ‘Y’ has a vertical line going between the two branches at the top) is unexplained.

\AM\ has two half-stanzas, clearly marked by the speaker notes in the margins.  The first one is by Weden: \emph{S. þ. e. xi. allir eins hæriar oðins tvnum i hꜹgguaz hværian dag.}  The second by Webthrithner: \emph{Val þæir kiosa ok riða vigi fra sitia mæirr vm sattir saman.}  This situation is also corrupt.

In order to restore st.~40 to a full four lines, two conjectural emendations have been adopted in the pres. ed.: In lines 1b--2 the refrain \emph{ef þïtt ǿði dugir ok þú Vaf·þrúðnir vitir} ‘if thy wisdom avails, and thou, Webthrithner, mightst know’ has been inserted from sts.~20 and 22, which also have ordinal numbers alliterating with vowels.  In line 3a the word \emph{allir} ‘all’ has been inserted from st.~41 to get vowel-alliteration with 3b \emph{ýtar}.}%
\speakernote{Óðinn kvað:}%
\llap{„}\edtext{\alst{S}ęg þat it}{\Afootnote{abbrev.~as \emph{‘S. þ. e.’} \AM}} \alst{ę}llipta, \hld\ \edtext{\emph{ef þïtt \alst{ǿ}ði dugir} &
\ind \emph{ok þú \alst{V}af·þrúðnir \alst{v}itir,}}{\lemma{ef \dots\ vitir,}\Afootnote{emend.; om.~all.}} &
hvar \edtext{\emph{\alst{a}llir}}{\Afootnote{emend.; om.~all.}} \hld\ \alst{ý}tar tu̇num ï &
\ind \alst{h}ǫggva⸗sk \alst{h}vęrjan dag.“\eva

\bvb\speakernoteb{Quoth Weden:}%
\emph{“Tell this} eleventh, \emph{if thy wisdom avails,} \\
\ind \emph{and thou, Webthrithner, mightst know} \\
where all men in yards \\
\ind strike each other every day.”\evb\evg


\bvg\bva\mssnote{\AM~3r/28, \RegiusProse~10r/17, \\
\Trajectinus~TODO, \Wormianus~TODO, \\
\Upsaliensis~TODO}%
\speakernote{Vaf·þrúðnir kvað:}
\Ballnote{This st.~is cited in \Gylfaginning~41:5, introduced by the words \emph{svá sem hér segir} ‘just as it says here’.  The immediately preceding prose passage \Gylfaginning~41:1--4 provides context: \emph{Þá mę́lti Gangleri: „All-mikill mann-fjǫlði er í Val-hǫll.  Svá njóta trú minnar, at all-mikill hǫfðingi er Óðinn, er hann stýrir svá miklum her.  Eða hvat er skemmtun Ein-herja, þá er þeir drekka eigi?“ Hárr segir: „Hvern dag, þá er þeir hafa klę́ðt⸗sk, þá her-vę́ða þeir sik ok ganga út í garð’inn ok berja⸗sk, ok fellir hverr annan.  Þat er leikr þeira.  Ok er líðr at dǫgurðar-máli, þá ríða þeir heim til Val-hallar ok setja⸗sk til drykkju.} ‘Then spoke Gangler: “There is an all-great gathering of men in Walhall.  Such makes me believe that Weden is an all-great chieftain when he wields such a great host.  But what is the entertainment of the Oneharriers when they do not drink?” High says: “Every day when they have got dressed then they gear themselves up for war and go out into the yard and fight, and each one slays the other.  That is their game.  And when it is getting close to supper-time then they ride home to Walhall and seat themselves at their drinking.”’

For the Oneharriers cf.~\textlink{Grimnismal}[18]/4, 24/3, 37/6, 53/3.}%
\llap{„}\alst{A}llir \edtext{\alst{ęi}n-hęrjar}{\Afootnote{so \GylfMS\emph{;} \emph{ęins hęrjar} \AM}} \hld\ \alst{Ó}ðins tu̇num ï &
\ind \alst{h}ǫggva⸗sk \alst{h}vęrjan dag; &
\edtrans{\alst{v}al þęir kjósa}{The slain they choose}{\Bfootnote{It is from this verbal phrase that the female agent noun \emph{val-kyrja} ‘\inx[P]{walkirries}[walkirrie]’ is derived.}} \hld\ ok ríða \alst{v}ígi frȧ, &
\ind \alst{s}itja męirr of \alst{s}áttir \alst{s}aman.“\eva

\bvb\speakernoteb{Quoth [Webthrithner]:}%
“All the \inx[P]{Oneharriers} in Weden’s yards \\
\ind strike each other every day. \\
The slain they choose and ride from the fray, \\
\ind then sit agreed once more together.”\evb\evg


\bvg\bva\speakernote{[Óðinn] kvað:}%
\mssnote{\Regius~8v/6, \AM~3v/1}%
\llap{„}Sęg þat it \alst{t}olpta, \hld\ hví þú \alst{t}íva rǫk &
\ind ǫll \alst{V}af·þrúðnir \alst{v}itir? &
Frȧ \alst{jǫ}tna ru̇num \hld\ ok \alst{a}llra goða &
\ind þú it \alst{s}annasta \alst{s}ęgir, &
\ind inn \alst{a}l-svinni \alst{jǫ}tunn.“\eva

\bvb\speakernoteb{Quoth [Weden]:}%
“Tell this twelfth, why thou the rakes of the Tews \\
\ind all, Webthrithner, shouldst know? \\
From the \inx[C]{rune}[runes] of the ettins and all the gods \\
\ind dost thou speak the truest, \\
\ind O all-wise ettin!”\evb\evg


\bvg\bva\mssnote{\Regius~8v/8, \AM~3v/2}%
\speakernote{[Vaf·þrúðnir] kvað:}%
\llap{„}Frȧ \alst{jǫ}tna ru̇num \hld\ ok \alst{a}llra goða &
\ind ek kann \alst{s}ęgja \alst{s}att, &
\ind því’t \alst{h}vęrn hęf’k \alst{h}ęim of komit, &
\edtext{\alst{n}íu kom’k hęima \hld\ fyr \alst{n}ifl-hęl neðan; &
\ind \alst{h}inig dęyja ór \alst{h}ęlju \alst{h}alir.}{\lemma{níu \dots\ halir. ‘Into nine \dots\ of Hell.’}\Bfootnote{Apparently lower Underworlds.  Cf.~Snorre, who writes about Lock’s daughter Hell in \Gylfaginning~34:6: \emph{Hel kastaði hann í Nifl-heim ok gaf henni vald yfir níu heimum, at hón skyldi skipta ǫllum vistum með þeim, er til hennar vǫ́ru sendir, en þat eru sótt-dauðir menn ok elli-dauðir.} ‘He [Weden, the Allfather] threw Hell into Nivelham and gave her dominion over nine Homes that she would deal out all provisions between those who were sent to her---but those are men dead of sickness and dead of old age.’  \textcite{FinnurEdda} considers \emph{ór hęlju} ‘out of Hell’ a later interpolation, probably for metrical reasons.}}“\eva

\bvb\speakernoteb{Quoth [Webthrithner]:}%
“From the runes of the ettins and all the gods \\
\ind I can speak truly, \\
\ind for I have come into each \inx[C]{Home}. \\
Into nine Homes I came beneath \inx[P]{Nivelhell}; \\
\ind that way men die out of \inx[P]{Hell}.”\evb\evg

\sectionline

\bvg\bva\mssnote{\Regius~8v/11, \AM~3v/4}%
\speakernote{Óðinn kvað:}%
\Ballnote{This question (Weden’s thirteenth) begins the third and final group of questions.  The subject shifts toward eschatology and the first two lines of each question-stanza feature the refrain found in st.~3/1--2.  For the events of the End Times, cf.~\textlink{Voluspa}[39]--55 and \Gylfaginning~51; for the world to come, cf.~\textlink{Voluspa}[56]--62, \Gylfaginning~52--53.}%
\llap{„}\alst{F}jǫlð ek \alst{f}ór, \hld\ \alst{f}jǫlð \alst{f}ręistaða’k, &
\ind \edtext{fjǫlð}{\Afootnote{abbrev.~as \emph{‘f.’} \Regius\AM}} \edtext{of}{\Afootnote{so \AM \emph{;} om.~\Regius}} \alst{r}ęynda’k \alst{r}ęgin; &
hvat lifir \alst{m}anna, \hld\ þȧ’s \edtext{inn \alst{m}ę́ra líðr &
\ind \alst{f}imbul-vetr}{\lemma{inn mę́ra \dots\ fimbul-vetr ‘the famed Fimble-winter’}\Bfootnote{The great winter of the end times.  The Fimblewinter is not mentioned in \textlink{Voluspa}’s description of the End Times, although bad weather and a blackening of the sun is mentioned in \textlink{Voluspa}[40]/3--4; snow during the End Times is also mentioned in \textlink{Hyndluljod}[41]/3.

\Gylfaginning~51:1--4 gives us a more detailed description of the Fimblewinter: \emph{Þá mę́lti Gangleri: „Hver tíðendi eru at segja frá um ragna-røkr? Þess hef ek eigi fyrr heyrt getit.“ Hárr segir: „Mikil tíðendi eru þaðan at segja ok mǫrg, þau in fyrstu, at vetr sá kemr, er kallaðr er fimbul-vetr.  Þá drífr snę́r ór ǫllum áttum.  Frost eru þá mikil ok vindar hvassir.  Ekki nýtr sólar.  Þeir vetr fara þrír saman ok ekki sumar milli.} ‘Then spoke Gangler: “Which tidings are to be told about the Twilight of the Reins? I have not heard it spoken of before.”  High says: “There are great tidings to be told thereof and many are they; the first that the winter will come which is called the Fimblewinter.  Then snow comes down from all directions.  The cold is great then and the winds biting.  The sun is of no use.  Those winters pass three together and no summer between them.’
For the immediately following prose in \Gylfaginning~51:5--6 see \textlink{Voluspa}[44] n.

Comparatively, an eschatological winter is known also from Zoroastrian Iranian sources, as described by \textcite{Hultgård2007}.  In one version the great winter sees numerous men and cattle killed, with the survivors finding refuge in a shelter.  Although the parallels to the Norse Fimblewinter are not great in number, the Iranian narratives are similar enough for \textcite[116--118]{Hultgård2007} to make note that they may share a common Indo-European origin with the Norse.

In recent scholarship \textcite{Gräslund2008} has identified the Fimblewinter as a traumatic reflex of the archeologically documented volcanic winter of 536--537 \ce, where a cloud of volcanic ash blocked out the light of the sun during the summer.  Gräslund rejects Hultgård’s comparative explanation and argues that the details of the Fimblewinter (the length of three years, the lack of summer, and the high mortality rate) are “highly precise” and therefore unlikely to stem from a long-lasting oral traditions \parencite[94--95]{Gräslund2008}.  This is undoubtedly a subjective line of argument, since the details are in fact rather sparse and not in agreement with the archeological record: all humans were not wiped out in 536, and the sun was not totally blocked out for three whole years.  Gräslund also conveniently omits the prediction of three winters of violence immediately preceding the three winters of the Fimblewinter as found in \Gylfaginning~51:5, since that can hardly be tied to his view of a historical reflex (unless we are to believe that there was a huge war in 532--534!)

As further support for his idea of historical trauma, \textcite[99--101]{Gräslund2008} cites an episode from the \emph{Kalevala} (songs 47--49) wherein the evil witch \emph{Louhi} steals the sun, moon, and fire and hides them in the Underworld.  According to Gräslund this would also be a reflex of the climate catastrophe of 536 \ce, but by making this connection he in fact undermines his own argument.  Gräslund is undoubtedly correct that the core concern of the Fimblewinter is not that the winter itself is harsh, but rather that the summer itself does not come, but such fears about malevolent supernatural interruptions to vital cosmic cycles are by no means isolated to sources composed after 536 \ce.
As early as \Rigveda\ (ca.~1500--1200 \bce) there is the myth of the Dawns imprisoned in a cave by the enemies of the Gods, thereby preventing the natural cycle between day and night (\cite[70--71]{JamBrer2020}), and the motif of Ettins trying to steal away things vital for the survival of mankind and the gods, not just the sun and moon, but also Thunder’s hammer, is found elsewhere in Norse myth (cf.~\textlink{Voluspa}[24]/4 n., Introduction to \textlink{Thrymskvida}).

It is clear that the return of the better side of the natural cycles (summer from winter, day from night, the full moon from the new) was of great concern to the ancients.  Though the volcanic winter of 536--537 was exceptionally severe, it cannot be (as Gräslund claims) the historical basis of every reflex of that motif when the examples demonstrably go back as far as the Bronze Age and likely beyond.  What is more, such worries could surely be fueled not just by volcanic trauma but by eclipses and cold summers (cf.~Eyv Lv 12--14 in \Skp\ 1) which have been occurring since time immemorial.  To pre-modern farmers, for whom the performance of the fields meant life and death, the return of summer and daylight were by no means felt inevitable, but were rather fraught with worry and therefore heralded with great rejoicing.
To conclude, I agree with Hultgård that the Fimblewinter may well be Indo-European.}} með \alst{f}irum?“\eva

\bvb\speakernoteb{Quoth Weden:}%
\emph{“Much I journeyed, much I tried, \\
\ind much I tested the Reins!} \\
What remains of men when the famed \inx[P]{Fimble-winter} \\
\ind passes amidst the folk?”\evb\evg


\bvg\bva\mssnote{\Regius~8v/13, \AM~3v/6, \\
\RegiusProse~17v/17, \Trajectinus~TODO, \\
\Wormianus~TODO, \Upsaliensis~TODO}%
\Ballnote{The coming ancestors of mankind will, like Ash and Emble (\textlink{Voluspa}[16]--17), emerge from wood.
This st.~is cited in \Gylfaginning, immediately following \Gylfaginning~53:11--13: \emph{En þar, sem heitir Hodd·mímis-holt, leyna⸗sk menn tveir í Surta-loga, er svá heita, Líf ok Leif·þrasir, ok hafa morgin-dǫggvar fyrir mat.  En af þessum mǫnnum kemr svá mikil kyn-slóð, at byggvi⸗sk heimr allr.  Svá sem hér segir:} ‘But there, at the place called Hoardmimer’s wood, two humans hide during the Flame of Surt, which are called Life and Lafthrasher, and they have morning dew for food.  But from these humans comes such a great line of descendants that the whole Home is settled.  Just as it says here.’}%
\speakernote{Vaf·þrúðnir kvað:}%
\llap{„}\alst{L}íf ok \edtext{\alst{L}ęif·þrasir}{\Afootnote{so \AM\RegiusProse\Wormianus; \emph{Líf·þrasir} \Regius\emph{;} \emph{Lęið·þrasir} \Trajectinus\emph{;} \emph{†lifþræser†} \Upsaliensis}}, \hld\ \edtext{ęn þau \alst{l}ęyna⸗sk munu}{\Afootnote{\emph{er þar leynazt meyjar} ‘who there hide, as maidens’ \Upsaliensis}} &
\ind ï \edtrans{\edtext{\edtext{\alst{h}olti}{\Afootnote{\emph{holdi} \AM}} \alst{H}odd·mïmis}{\Afootnote{\emph{Mïmis holdi} ‘Mimer’s corpse’ \Upsaliensis}}}{in Hoardmimer’s wood}{\Bfootnote{Perhaps the hollowed-out trunk of \inx[P]{Ugdrassle’s Ash.}}}; &
\edtext{\alst{m}orgin-dǫggvar}{\Afootnote{\emph{morgin-dǫggva} \Upsaliensis}} \hld\ \edtext{þau sér at \alst{m}at hafa}{\Afootnote{\emph{þę́r} \Upsaliensis}}; &
\ind \edtext{\edtext{ęn}{\Afootnote{om. \Regius}} þaðan af \alst{a}ldir}{\Afootnote{\emph{ok þar of aldr} \Upsaliensis}} \edtext{\alst{a}la⸗sk}{\Afootnote{\emph{allar} \Trajectinus}}.“\eva

\bvb\speakernoteb{Quoth Webthrithner:}%
“\inx[P]{Life} and \inx[P]{Lafthrasher}---but they will hide themselves \\
\ind in \inx[P]{Hoardmimer}’s wood. \\
Morning dew will they have as food, \\
\ind but thereof is mankind begotten.”\evb\evg


\bvg\bva\mssnote{\Regius~8v/15, \AM~3v/8}%
\speakernote{Óðinn kvað:}%
\llap{„}\alst{F}jǫlð \edtext{ek \edtext{\alst{f}ór, \hld\ \alst{f}jǫlð \alst{f}ręistaða’k, &
\ind fjǫlð of \alst{r}ęynda’k \alst{r}ęgin}{\lemma{fór \dots\ ręgin}\Afootnote{abbrev.~as \emph{‘f. f. e. f.’} \Regius}}}{\lemma{ek \dots\ ręgin}\Afootnote{abbrev.~as \emph{‘e. f. f. e. f.’} \AM}}; &
hvaðan kømr \alst{s}ól \hld\ ȧ inn \alst{s}létta himin, &
\ind \edtext{þȧ’s}{\Afootnote{so (\emph{‘þa ær’}) \AM \emph{;} \emph{‘þa þa er’} \Regius}} \edtrans{þęssa}{this one}{\Bfootnote{The current sun, as explained in the following st.}} hęfr \edtrans{\alst{F}ęnrir}{Fenrer}{\Bfootnote{It is not certain that this is the same “Fenrerswolf” which fights against Weden (cf.~st.~53 below);
the word \emph{fęnrir}, which prob.~originally means ‘fen-dweller’, may here simply be a generic poetic synonym for ‘wolf’.  For the wolves which chase and eventually devour the sun and moon see \textlink{Voluspa}[39]--40 (n.), \textlink{Grimnismal}[40], \Gylfaginning~12.}} \alst{f}arit?\eva

\bvb\speakernoteb{Quoth Weden:}%
\emph{“Much I journeyed, much I tried, \\
\ind much I tested the Reins!} \\
Whence comes the Sun onto the smooth heaven \\
\ind when \inx[P]{Fenrer} has this one killed?”\evb\evg


\bvg\bva\mssnote{\Regius~8v/16, \AM~3v/9, \\
\RegiusProse~17v/20, \Trajectinus~TODO, \\
\Wormianus~TODO, \Upsaliensis~TODO}%
\Ballnote{This st.~is cited in \Gylfaginning, immediately following \Gylfaginning~53:14--15: \emph{Ok hitt mun þér undar⸗ligt þykkja er sól’in hefir getit dóttur eigi ó·fegri en hón er, ok ferr sú þá stigu móður sinnar.  Svá sem hér segir:} ‘And it will seem wondrous to thee that the Sun has begotten a daughter no less fairer than she herself is, and that one then journeys upon her mother’s paths.  Just as it says here:’
The immediately following passage \Gylfaginning~53:16--17 concludes the account of the world to come and the mythological section of \Gylfaginning\ as a whole: \emph{En nú ef þú kannt lengra fram at spyrja, þá veit ek eigi, hvaðan þér kemr þat, fyrir því at engan mann heyrða ek lengra segja fram aldar-far’it, ok njót⸗tu nú sem þú namt.} ‘But now, if thou canst ask further ahead then I do not know for what reason thou mightst do so, for I have heard no man tell further ahead about the progression of the world; and now enjoy what thou hast learned.’}%
\speakernote{Vaf·þrúðnir kvað:}%
\llap{„}\alst{Ęi}na dóttur \hld\ berr \edtrans{\alst{a}lf-rǫðull}{Elf-wheel}{\Bfootnote{A rare poetic synonym (\emph{hęiti}) for the sun.  It occurs independently in two other places: the Eddic \textlink{Skirnismal}[4]/3a and a Scaldic loose stanza by Eanwind Finnson (Eyv Lv 9 in \Skp\ 1).  It also appears in two lists of names for the sun: \Skaldskaparmal\ 69 and Þul \emph{Sólar} 1/7 (in \Skp\ 3), but these do not count as independent attestations since they are antiquarian sources drawing on earlier poetry.

Depending on the age of the cpd., \emph{alf-} may either reflect the semantics of PIE \emph{albʰós} ‘white’ (whence Latin \emph{albus} ‘id.’) or ON \emph{alfr} ‘elf’, the latter word deriving from the former and originally meaning ‘white one’.  That is not to say that \emph{alf-rǫðull} goes back to PIE, but only that the sense “white” must have existed at some point in Germanic prehistory.
The second element, \emph{rǫðull}, is not the normal ON word for ‘wheel’ (\emph{hvél} \char`~\ \emph{hjól}).  It is inherited from PGmc.~*\emph{radulaz} \char`~\ *\emph{raduraz} (> OE \emph{rǫdor} ‘heaven, sky’, OS \emph{radul, radur} ‘id.’), from *\emph{raþą} ‘wheel’ (> German \emph{Rad} ‘id.’) + the nominal suffix *\emph{-ulaz} \char`~\ *\emph{-uraz} ‘(habitually) doing’ and thus means something like ‘circler, turner, revolver’.
The PIE root of *\emph{raþą} is *\emph{Hreth₂-} (whence also Latin \emph{rota} ‘wheel’, Sanskrit \emph{rátʰa} ‘chariot’; \cite[405]{Kroonen2013}).}}, &
%NOTE to self: Kroonen2013 has no entries for *albiz and *raduraz!
\ind áðr \edtext{hana}{\Afootnote{\emph{hęnni} \Upsaliensis}} \alst{F}ęnrir \alst{f}ari; &
\edtext{sú skal \alst{r}íða, \hld\ þȧ’s \alst{r}ęgin dęyja}{\Afootnote{\emph{sú mun renna \hld\ eða ríða} (unmetrical and wo.~doubt corrupt) \Upsaliensis}}, &
\ind \edtext{\alst{m}óður}{\Afootnote{\emph{ręgin} \Upsaliensis}} brautir \alst{m}ę́r.“\eva

\bvb\speakernoteb{Quoth Webthrithner:}%
“A single daughter the Elf-wheel \ken{sun} bears \\
\ind before Fenrer might kill her. \\
She shall ride when the Reins die \\
\ind her mother’s roads, a maiden.”\evb\evg


\bvg\bva\mssnote{\Regius~8v/18, \AM~3v/10}%
\speakernote{Óðinn kvað:}%
\edtext{\alst{F}jǫlð ek \edtext{\alst{f}ór, \hld\ \alst{f}jǫlð \alst{f}ręistaða’k, &
\ind fjǫlð of \alst{r}ęynda’k \alst{r}ęgin}{\lemma{fór \dots\ ręgin}\Afootnote{abbrev.~as \emph{‘f. f.’} \Regius}}}{\lemma{Fjǫlð \dots\ ręgin}\Afootnote{abbrev.~as \emph{‘F. e. f. f. e. f.’} \AM}}; &
hvęrjar ’ro \alst{m}ęyjar, \hld\ es líða \alst{m}ar yfir, &
\ind \alst{f}róð-gęðjaðar \alst{f}ara?\eva

\bvb\speakernoteb{Quoth Weden:}%
\emph{“Much I journeyed, much I tried, \\
\ind much I tested the Reins!} \\
Who are the maidens that pass o’er the sea? \\
\ind Wise-minded they go.”\evb\evg


\bvg\bva\mssnote{\Regius~8v/19, \AM~3v/11}%
\Ballnote{The identity of these maidens is entirely obscure.  Considering all other questions introduced with the refrain \emph{Fjǫlð ek fór} \dots\ have something to do with the end times, they ought to as well, but that hardly clarifies their function.  They are probably to be identified with the maidens about which Weden asks in \textlink{Baldrsdraumar}[12].}%
\speakernote{Vaf·þrúðnir kvað:}%
\alst{Þ}ríar \alst{þ}jóð-á\emph{a}r \hld\ falla \alst{þ}orp yfir &
\ind \alst{m}ęyja \alst{M}ǫg·þrasis; &
\alst{h}amingjur ęinar \hld\ \edtext{þę́r’s}{\Afootnote{so \emph{(‘þær ær’)} \AM \emph{;} \emph{þęira} \Regius}} ï \alst{h}ęimi eru, &
\ind þó þę́r með \alst{jǫ}tnum \alst{a}la⸗sk.\eva

\bvb\speakernoteb{Quoth Webthrithner:}%
“Three great rivers fall over the house \\
\ind of the maidens of Maythrasher. \\
They are the only Hamings in the Home, \\
\ind although they are raised amidst ettins.”\evb\evg


\bvg\bva\mssnote{\Regius~8v/21, \AM~3v/13}%
\speakernote{Óðinn kvað:}%
\llap{„}\edtext{\alst{F}jǫlð ek \edtext{\alst{f}ór, \hld\ \alst{f}jǫlð \alst{f}ręistaða’k, &
\ind fjǫlð of \alst{r}ęynda’k \alst{r}ęgin}{\lemma{fór \dots\ ręgin}\Afootnote{abbrev.~as \emph{‘f. f.’} \Regius}}}{\lemma{Fjǫlð \dots\ ręgin}\Afootnote{abbrev.~as \emph{‘F. e. f. f. e. f.’} \AM}}; &
hvęrir ráða \alst{ę́}sir \hld\ \alst{ęi}gnum goða, &
\ind þȧ’s \alst{s}loknar \edtrans{\alst{S}urta-logi}{the flame of Surt}{\Bfootnote{The flame which reaches up to Heaven itself and scorches the whole world; see \textlink{Voluspa}[50], 54.}}?“\eva

\bvb\speakernoteb{Quoth Weden:}%
\emph{“Much I journeyed, much I tried, \\
\ind much I tested the Reins!} \\
Which Eese will rule the estates of the Gods \\
\ind when the flame of \inx[P]{Surt} goes out?”\evb\evg


\bvg\bva\mssnote{\Regius~8v/22, \AM~3v/14, \\
\RegiusProse~17v/12, \Trajectinus~TODO, \\
\Wormianus~TODO, \Upsaliensis~TODO}%
\Ballnote{This st.~is cited in \Gylfaginning~53:10, introduced by the words \emph{Svá er sagt} ‘So it is said’.  For the preceding prose passage \Gylfaginning~53:1--9 see \textlink{Voluspa}[56] n.}%
\speakernote{Vaf·þrúðnir kvað:}%
\llap{„}\alst{V}íð·arr ok \alst{V}ȧli \hld\ byggva \alst{v}é goða, &
\ind þȧ’s \alst{s}loknar \alst{S}urta-logi; &
\edtext{\alst{M}óði}{\Afootnote{\emph{†megi†} \RegiusProse}} ok \alst{M}agni \hld\ skulu \alst{M}jǫllni hafa &
\ind \edtrans{\edtext{\alst{V}ingnis}{\Afootnote{\emph{ok vinna} \Regius\emph{;} add. \emph{synir} \Upsaliensis}} at \edtext{\alst{v}íg-þroti}{\Afootnote{\emph{†vigroþi†} \Upsaliensis}}}{after Wingner expires in war}{\Bfootnote{After Thunder dies in his fight against the \inx[P]{Middenyardswyrm}, for which see \textlink{Voluspa}[53].}}.“\eva

\bvb\speakernoteb{Quoth Webthrithner:}%
“\inx[P]{Wider} and \inx[P]{Wonnel} will settle the \inx[C]{wigh}[wighs] of the Gods \\
\ind when the flame of Surt goes out. \\
\inx[P]{Mood} and \inx[P]{Main} shall have \inx[P]{Millner} \\
\ind after \inx[P]{Wingner} expires in war.”\evb\evg


\bvg\bva\mssnote{\Regius~8v/24, \AM~3v/16}%
\speakernote{Óðinn kvað:}%
\llap{„}\edtext{\alst{F}jǫlð ek \edtext{\alst{f}ór, \hld\ \alst{f}jǫlð \alst{f}ręistaða’k, &
\ind fjǫlð of \alst{r}ęynda’k \alst{r}ęgin}{\lemma{fór \dots\ ręgin}\Afootnote{abbrev.~as \emph{‘f.’} \Regius}}}{\lemma{Fjǫlð \dots\ ręgin}\Afootnote{abbrev.~as \emph{‘F. e. f.’} \AM}}; &
hvat verðr \alst{Ó}ðni \hld\ at \alst{a}ldr-lagi, &
\ind \edtrans{þȧ’s \edtext{of}{\Afootnote{so \AM \emph{;} om. \Regius}} \alst{r}júfa⸗sk \alst{r}ęgin?}{when the Reins are ripped?}{\Bfootnote{Formulaic; see note to \textlink{Baldrsdraumar}[14]/4.}}“\eva

\bvb\speakernoteb{Quoth Weden:}%
\emph{“Much I journeyed, much I tried, \\
\ind much I tested the Reins!} \\
What brings Weden’s life to an end \\
\ind when the Reins are ripped?”\evb\evg


\bvg\bva\mssnote{\Regius~8v/25, \AM~3v/17}%
\Ballnote{For Weden’s fight with the Wolf and Wider’s subsequent avenging see \textlink{Voluspa}[51]--52, \textlink{Grimnismal}[17], 24.}%
\speakernote{Vaf·þrúðnir kvað:}%
\llap{„}\alst{U}lfr glęypa \hld\ mun \alst{A}lda-fǫðr, &
\ind \edtrans{þęss mun \alst{V}íð·arr \edtext{\emph{\alst{v}}reka}{\Afootnote{metr. emend.; \emph{reka} \Regius}}}{that will Wider avenge}{\Bfootnote{The archaic initial \emph{vr-} in \emph{vreka} must be restored for metrical reasons \parencite[37]{Haukur2016}.  This serves as a solid dating criterion for composition before ca.~1000 \ce; cf.~\textlink{Havamal}[26]/2 n.}}; &
\alst{k}alda \alst{k}japta \hld\ hann \alst{k}lyfja mun &
\ind \edtrans{\alst{v}itnis}{the beast}{\Afootnote{so \AM\emph{;} \emph{Vingnis} ‘Wingner’ (prob. influenced by st.~51) \Regius}} \alst{v}ígi at.“\eva

\bvb\speakernoteb{Quoth Webthrithner:}%
“The Wolf will devour the Father of Mankind; \\
\ind that will Wider avenge. \\
The cold jaws he will split apart \\
\ind of the beast at the battle.”\evb\evg


\bvg\bva\speakernote{Óðinn kvað:}%
\mssnote{\Regius~8v/27, \AM~3v/19}%
\llap{„}\edtext{\alst{F}jǫlð ek \alst{f}ór, \hld\ \edtext{\alst{f}jǫlð \alst{f}ręistaða’k, &
\ind fjǫlð of \alst{r}ęynda’k \alst{r}ęgin}{\lemma{fjǫlð \dots\ ręgin}\Afootnote{abbrev.~as \emph{‘f.’} \Regius}}}{\lemma{Fjǫlð \dots\ ręgin}\Afootnote{abbrev.~as \emph{‘F. e. f.’} \AM}}; &
\edtext{hvat mę́lti Óðinn, \hld\ áðr \edtrans{ȧ bál \edtext{stigi}{\lemma{stigi,}\Afootnote{add.~\emph{ok} \AM}}}{step onto the pyre}{\Bfootnote{The phrase \emph{stíga ȧ} ‘step onto, mount’ is also used to refer to one stepping aboard a ship or mounting a horse (see \CV: \emph{stíga} for citations).
Its use for a person being borne onto the funeral pyre has been compared with \Beowulf\ 1118b: \emph{gu̇ð-rinc á·stâh} ‘the war-champion mounted [his pyre]’, but the interpretation of that line is controversial.  \textcite[186]{KlaeberBeowulf} instead follow Grundtvig in emending \emph{gu̇ð-rinc} to \emph{gu̇ð-réc} ‘war-smoke’ and compare it with \Beowulf\ 3144b (\emph{wudu-réc á·stâh} ‘wood-smoke rose up’, which also describes a cremation.  Thus, the present stanza “almost certainly refers not to Baldr but to Óðinn, probably imagined to mount the pyre in order to set fire to it” \parencite[186]{KlaeberBeowulf}.}}, &
\ind \alst{s}jalfr ï ęyra \alst{s}yni?}{\lemma{Hvat mę́lti Óðinn \dots\ ï ęyra syni? ‘What said Weden \dots\ into the ear of his son?’}\Bfootnote{Weden poses the eighteenth, unknowable, question.  Only he himself can know what he whispered in Balder’s ear, and thus he has won the contest.  For Balder’s death cf.~notes to \textlink{Voluspa}[31]--33 and \textlink{Baldrsdraumar}[8]--11.
Weden asks the same question in Heathric’s Riddles in \HervararSaga, as the 37th and final riddle (for critical discussion cf.~Heiðr 84/3 n.~in \Skp\ 8.).  The half-stanza there is clearly related to \Vafthrudnismal\ 54: \emph{Sęg þat þá fyrst, ef þú ert hvęrjum konungi vitrari: // Hvat mę́lti \alst{Ó}ðinn \hld\ í \alst{ęy}ra Baldri, / áðr hann vę́ri á \alst{b}ál of \alst{b}orinn?} ‘Then tell this first if thou art cleverer than every king: // What spoke Weden in Balder’s ear / before he would be borne onto the pyre?’}}“\eva

\bvb\speakernoteb{Quoth Weden:}%
\emph{“Much I journeyed, much I tried, \\
\ind much I tested the Reins!} \\
What spoke Weden before he would step onto the pyre, \\
\ind himself, into the ear of his son?”\evb\evg


\bvg\bva\speakernote{Vaf·þrúðnir kvað:}%
\mssnote{\Regius~8v/28, \AM~3v/19--22}%
\llap{„}\alst{Ęy} \edtext{mann-\emph{g}i}{\Afootnote{emend.; \emph{manni} (dat.~sg., impossible since a subject to \emph{vęit} is needed) \Regius\AM}} vęit, \hld\ hvat þú ï \edtrans{\alst{á}r-daga}{in days of yore}{\Bfootnote{Relative to the events of the poem, Balder’s death}} &
\ind \alst{s}agðir ï ęyra \alst{s}yni; &
\edtrans{\alst{f}ęigum}{fey}{\Bfootnote{“Death-doomed”, a word with fatalistic connections.  Webthrithner realises that he was bound to die from the moment he proposed the wager (st.~19), as no being can outwit the god of wisdom.}} munni \hld\ \edtext{mę́lta’k}{\Afootnote{add.~\emph{umb} (\emph{‘um’}) \AM}} \edtrans{mïna \alst{f}orna stafi}{my ancient staves}{\Bfootnote{Referencing st.~1.}} &
\ind ok of \alst{r}agna \alst{r}ǫk; &
nú við \alst{Ó}ðin \hld\ dęilda’k mïna \edtrans{\alst{o}rð-spęki}{word-wisdom}{\Bfootnote{Referencing st.~5.}}; &
\ind þú est ę́ \alst{v}ísastr \edtrans{\alst{v}era}{of men}{\Bfootnote{\emph{verr} means ‘husband, man’ and is here used for reasons of alliteration; it does not imply that Weden is not a God.}}.“\eva

\bvb\speakernoteb{Quoth Webthrithner:}%
“Never will man know what thou in days of yore \\
\ind spokest in the ear of thy son. \\
With a \inx[C]{fey} mouth I spoke my ancient \inx[C]{stave}[staves] \\
\ind and about the Rakes of the Reins. \\
Now with Weden have I shared my word-wisdom--- \\
\ind thou art for ever wisest of men!”\evb\evg

\sectionline
